% ===========================================================
% 《线性代数9讲》第 9 讲  二次型 —— 片段 A（本讲开头）
% 源：线代9讲强化.pdf  PDF p104–p112（印刷页 93–101），共 9 页
%
% 处理说明：
%   1) PDF p104（印刷页 93）页首「三向解题法」思维导图（位于讲标题「第9讲 二次型」
%      之下、正文「一、$f=x^{\mathrm T}Ax$ 中 $A$ 的表示」之上）按 AGENTS.md §7.1 决定 2
%      **不转写、不重绘**，只留一行 `% FIGURE-OPD: 93 三向解题法思维导图（按规则不保留）`；
%      其上的蓝色小标题「三向解题法」随之删除（留痕见下）。
%   2) OPD 方法论标记剔除（依 AGENTS.md §7.1 决定 2 与本片要求）：
%      · \fsec 标题里只删 OPD 括号，**保留方法名**：
%        「一、$f=x^{\mathrm T}Ax$ 中 $A$ 的表示」（$O_1$（盯住目标 1）$+D_1$（常规操作）
%          $+D_{23}$（化归经典形式））
%        「二、配方法与正交变换法的异同」（$O_2$（盯住目标 2）$+D_{23}$（化归经典形式））
%        「三、伪配方法」（$O_3$（盯住目标 3）$+D_{23}$（化归经典形式））
%        「四、正交变换法的传递性」（$O_4$（盯住目标 4）$+D_1$（常规操作）
%          $+D_{23}$（化归经典形式））
%        「五、合同的判定与手段」（$O_5$（盯住目标 5）$+D_1$（常规操作）
%          $+D_{23}$（化归经典形式））
%      · 正文蓝色纯 OPD 代码（带箭头旁注）与蓝色方法论提示语一律以
%        `% OPD 蓝色标注（原稿）：……` 注释留痕，不排入正文：
%        PDF p104（93）节标题上方蓝色小标签「小题考法」
%        PDF p105（94）式下蓝色提示语「目的：凑二次型定义」、「按 $f=x^{\mathrm T}Ax$ 的方向走」
%          及其 OPD 码「$D_{23}$（化归经典形式）」
%        PDF p106（95）四处「→$D_{22}$（转换等价表述）」及蓝色提示语
%          「要掌握这两个等价的语言表达并能相互转化」（2 处）、
%          「两种方法一定要搞清，不要只会套公式」
%        PDF p107（96）「→$D_{22}$（转换等价表述）」
%        PDF p108（97）「$D_{23}$（化归经典形式）」及其蓝色提示语
%          「没有平方项，创造平方项」
%        PDF p110（99）例 9.7 题干右侧「→基本局面」（原稿蓝色）
%        PDF p111（100）「$D_{23}$（化归经典形式）」
%        PDF p112（101）例 9.9 题干右侧「→变换手段」
%      · 页边竖排模块名「形式化归体系块」（PDF p104、p106 页右缘，PDF p109 页左缘）
%        为 OPD 体系块标记，删除并留痕。
%      · **数学符号一律照抄**：$A,\ x,\ \Lambda,\ \lambda,\ p_A,\ q_A$ 等不受影响。
%   3) **数学操作旁注照原文排入**（`\quad(\text{...})` 内联或行内）：包括
%      「主对角线元素不变」「利用内积定义 $(\alpha,\beta)=(\beta,\alpha)$」
%      「$1$ 和 $\frac34$ 不一定是特征值」「$(-1)$ 倍加至」
%      「$C$ 不一定由 $A$ 的特征向量拼成，$\Lambda$ 也不一定是由 $A$ 的特征值拼成的」
%      「实对称矩阵不同特征值对应的特征向量正交，故只需单位化」
%      「$B$ 的特征值为 $k,6,0$，$|B|=k\cdot6\cdot0=0$，$A,B$ 合同，故 $A$ 有一个特征值 $0$」
%      「因为其行列式为 $0$，所以此矩阵不可逆……把整个问题的性质改变了」
%      「配方时，不一定要先配 $x_1$……就先配哪个」「盯着 $A$ 的对角线元素，平方项分别提出 2，3，1」。
%      原稿的笔误、错别字、编号一律照抄。
%   4) ⚠️ 文本模式关系符一律写 `$\ne$`/`$\le$`/`$\ge$`；$\fsec$ 标题内不用 `\cdots`。
%   5) **本 9 页无「习题」栏**（已逐页核实）：全部为知识点正文 + 例 9.1～例 9.9 的例题与解答。
%   6) 本片无位图插图（各页右上角蓝色二维码为装饰，未转写）；无 `fdeep`（9 讲正文不使用）。
%   7) 版式还原说明：PDF p106（95）原稿把「①配方法」「②正交变换法」两条并排收在一个左大括号内
%      （左侧竖排提示语即第 2) 条所列已删提示语），单栏宽度无法并排，故按原编号顺序改为
%      顺序段落，先①后②；同页 $f(x)=x^{\mathrm T}Ax=$ 以 \fmll 单独成行，随后接①②两条。
%      PDF p111（100）例 9.8 的两条页缘蓝色说明（「配方时……」与「盯着 $A$ 的对角线元素……」）
%      按数学操作旁注内联排入，位置由版心左缘/公式下方移到公式之后。
%
% 接缝说明：
%   · 本片首行 = 第 9 讲开篇（PDF p104），其上无第 8 讲残余内容；
%   · 末行 = 例 9.9「③令 $D_1x=D_2y$，求 $D_2^{-1}D_1$」的 $D$ 矩阵计算结果
%     （PDF p112 末，行末为逗号），该例的收尾结论「故 $A=D^{\mathrm T}BD$」在
%     PDF p113（印刷页 102），由后续片段承接；例 9.9 自 PDF p112 首行起，未跨片前结束。
%   · 例 9.1 题面在 PDF p104（93）末，其解答在 PDF p105（94）首，跨页但同片，已在同一张
%     `fcard` 内承接。
% ===========================================================

\fchap{第9章\quad 二次型}

% -----------------------------------------------------------
% PDF p104（印刷页 93）
% -----------------------------------------------------------

% FIGURE-OPD: 93 三向解题法思维导图（按规则不保留）

% OPD 蓝色标注（原稿）：页首蓝色小标题「三向解题法」（方法论板块，按规则不保留）
% OPD 蓝色标注（原稿）：节标题上方蓝色小标签「小题考法」（方法论标签，已删）
% OPD 蓝色标注（原稿）：（$O_1$（盯住目标 1）$+D_1$（常规操作）$+D_{23}$（化归经典形式））（\fsec 标题中的 OPD 括号已删，保留方法名）
% OPD 蓝色标注（原稿）：页右缘竖排模块名「形式化归体系块」（OPD 体系块标记，已删）

\fsec{深化一　二次型的矩阵表示与合同}

\begin{fdeep}{二次型的矩阵为什么必须取对称的}
二次型是二次齐次多项式，例如
\fml{f(x_1,x_2)=x_1^{2}+4x_1x_2+3x_2^{2}}
把它写成矩阵形式时，交叉项的系数要“平分”到两个位置上：
\fml{f=\begin{pmatrix}x_1&x_2\end{pmatrix}
\begin{pmatrix}1&2\\2&3\end{pmatrix}\begin{pmatrix}x_1\\x_2\end{pmatrix}=x^{\mathrm T}Ax}
\textbf{关键规定}：$A$ 必须取\textbf{对称矩阵}（$a_{ij}=a_{ji}=\frac12\times$ 交叉项系数）。
原因不是“习惯”，而是有必然性：
\fml{x^{\mathrm T}Ax=x^{\mathrm T}\Bigl(\frac{A+A^{\mathrm T}}{2}\Bigr)x}
（因为 $x^{\mathrm T}Ax$ 是 $1\times1$ 矩阵，等于自身的转置，故 $x^{\mathrm T}Ax=x^{\mathrm T}A^{\mathrm T}x$）。
也就是说\textbf{只有对称部分对二次型有贡献}，反对称部分被自动消掉。
于是“二次型 $\leftrightarrow$ 对称矩阵”是\textbf{一一对应}——每个二次型唯一确定一个对称矩阵，
反之亦然。这就是为什么讲二次型时永远先取对称矩阵。

\fbref{}：服务考纲〔二次型〕“掌握二次型及其矩阵表示”。
考场上最常见的失分是\textbf{把交叉项系数直接填进矩阵}（写成 $a_{12}=4$ 而不是 $2$），
或者填出一个不对称的矩阵。\must{记住}“系数平分、矩阵对称”即可避免，
而“反对称部分对二次型无贡献”（第 4 章已证）正是这条规定的理论根据。
\end{fdeep}

\begin{fdeep}{合同 vs 相似：两种变换、两套不变量}
二次型换到新变量下，矩阵的变化是\textbf{合同}：
\fml{x=Cy\ \Longrightarrow\ x^{\mathrm T}Ax=y^{\mathrm T}(C^{\mathrm T}AC)y,
\qquad B=C^{\mathrm T}AC}
把三种关系统一对照（这是全书最该背下的一张表）：

\begin{center}
\begin{tabular}{@{}llll@{}}
\toprule
关系 & 变换形式 & 要求 & 不变量 \\
\midrule
等价（相抵） & $B=PAQ$ & $P,Q$ 可逆 & 秩 \\
相似 & $B=P^{-1}AP$ & $P$ 可逆 & 特征值、迹、行列式、秩 \\
合同 & $B=C^{\mathrm T}AC$ & $C$ 可逆 & 正负惯性指数（秩 + 定性） \\
\bottomrule
\end{tabular}
\end{center}

\textbf{三者的强弱关系}：
\begin{itemize}[leftmargin=2.2em,itemsep=0.22em]
\item \textbf{相似 $\Rightarrow$ 等价}（取 $Q=P^{-1}$ 即可），但反之不成立；
\item \textbf{合同 $\Rightarrow$ 等价}（取 $P=C^{\mathrm T}$、$Q=C$），但反之不成立；
\item \textbf{相似与合同互不包含}：相似不保证对称性保持（$P^{-1}AP$ 未必对称），
      合同不保证特征值不变（$C^{\mathrm T}AC$ 的特征值一般会变）。
\end{itemize}
\textbf{唯一的交集}：若 $C$ 取\textbf{正交矩阵} $Q$（$Q^{\mathrm T}=Q^{-1}$），则
$Q^{\mathrm T}AQ=Q^{-1}AQ$ \textbf{既是合同又是相似}——这就是“正交对角化”的特殊地位：
它同时保住了惯性指数和特征值。

\fbref{}：服务考纲〔二次型〕“了解合同变换与合同矩阵的概念”，并与〔矩阵的特征值和特征向量〕的“相似”对照。
考场上判断题“A 与 B 是相似还是合同”时，按上表逐项比不变量：
\textbf{看特征值就是想相似，看正负惯性指数就是想合同}。
用错关系会导致整题方向错误——这是二次型部分最根本的一个辨析点。
\end{fdeep}

\begin{fdeep}{$A$ 与 $B$ 合同的判定：比正负惯性指数}
\textbf{核心判据}：$n$ 阶实对称矩阵 $A$ 与 $B$ \textbf{合同} $\iff$ 它们的正惯性指数相同
\textbf{且}秩相同；等价说法是“对应的二次型有相同的规范形”（两矩阵必合同于
$\operatorname{diag}(E_p,-E_{r-p},O)$，全部合同类由 $(p,\ r-p)$ 分类）。
考场三步：先比\textbf{秩}（秩不等直接不合同，因为合同必等价、等价保秩）；
秩相等再比\textbf{正惯性指数 $p$}（用顺序主子式法数正系数，或数正特征值个数）；
两项都相等则合同，否则不合同。
另有一条常考的唯一性结论：若 $A^{\mathrm T}=A,\ B^{\mathrm T}=B$ 且对一切 $x$ 有
\fml{x^{\mathrm T}Ax=x^{\mathrm T}Bx\ \Longrightarrow\ A=B}
（二次型与对称矩阵\textbf{一一对应}：$x_ix_j\ (i\ne j)$ 项系数的一半即 $a_{ij}$，$x_i^{2}$ 项系数即 $a_{ii}$。）

\fbref{}：服务考纲〔二次型〕“了解合同变换与合同矩阵”“了解…惯性定理”。选择题“下列矩阵中与 $A$ 合同的是”，标准解法就是\textbf{算秩 + 算正惯性指数}这两个数，根本不必求出变换矩阵。再加一句“对称矩阵由二次型唯一决定”，可秒杀一类“已知 $x^{\mathrm T}Ax=x^{\mathrm T}Bx$，证明 $A=B$”的证明题。
\end{fdeep}

\begin{fdeep}{二次型的秩：三种算法与 rank$(A^{\mathrm T}A)$}
二次型的\textbf{秩}定义为它矩阵 $A$ 的秩，也等于标准形（规范形）中\textbf{非零平方项的个数}
——合同变换保秩，而对角阵的秩就是非零对角元的个数。由此得三条实用结论：
（1）秩是合同不变量中最弱的一个：\textbf{秩不同，一定不等价、不合同、不相似}；
（2）若 $f=\sum_{i=1}^{s}\bigl(a_{i1}x_1+a_{i2}x_2+\dots+a_{in}x_n\bigr)^{2}$，
记系数矩阵为 $A_{s\times n}$，则 $f$ 的矩阵是 $A^{\mathrm T}A$，而
\fml{\operatorname{rank}(A^{\mathrm T}A)=\operatorname{rank}A}
故 $f$ 的秩 $=\operatorname{rank}A\le\min\{s,n\}$；
（3）特别地，$s<n$ 时 $f$ \textbf{不可能正定}（最多半正定）——正定要求有 $n$ 个平方项
且系数矩阵满秩。

\fbref{}：服务考纲〔二次型〕“了解二次型秩”。“$f$ 是若干个完全平方之和”是命题人最爱的给法：一眼写出矩阵 $A^{\mathrm T}A$，秩就等于系数矩阵的秩，比展开再求秩快得多；同时它解释了为什么“$n$ 元二次型写成少于 $n$ 个平方和”必然不正定——半正定与正定的分水岭就在这里。
\end{fdeep}

\fsec{一、$f=x^{\mathrm T}Ax$ 中 $A$ 的表示}

（1）给出非对称矩阵 $B$，令 $A=\dfrac{B+B^{\mathrm T}}{2}$，则 $A=A^{\mathrm T}$.

（2）通过题设或基本变形显化出 $A$.

\begin{fcard}{例9.1}
已知三元二次型表示为 $f(x)=x^{\mathrm T}\begin{bmatrix}1&1&0\\0&1&1\\0&0&1\end{bmatrix}x$，则 $f$ 的规范形为（$\quad$）.

（A）$y_1^2-y_2^2+y_3^2$\qquad（B）$y_1^2+y_2^2+y_3^2$

（C）$-y_1^2-y_2^2-y_3^2$\qquad（D）$y_1^2-y_2^2-y_3^2$

% ---- PDF p105（印刷页 94）----

【解】应选（B）.

二次型的矩阵为 $A=\begin{bmatrix}1&\frac12&0\\[2pt]\frac12&1&\frac12\\[2pt]0&\frac12&1\end{bmatrix}$ $\quad(\text{主对角线元素不变})$，其各阶顺序主子式为

\fml{\Delta_1=1>0,\qquad \Delta_2=\begin{vmatrix}1&\frac12\\[2pt]\frac12&1\end{vmatrix}=\frac34>0,\qquad \Delta_3=\begin{vmatrix}1&\frac12&0\\[2pt]\frac12&1&\frac12\\[2pt]0&\frac12&1\end{vmatrix}=\frac12>0.}

所以 $A$ 是正定矩阵，其正惯性指数 $p=3$，因此二次型的规范形为 $y_1^2+y_2^2+y_3^2$.
\end{fcard}

\begin{fcard}{例9.2}
已知
\fml{\alpha_1=\begin{bmatrix}1\\2\end{bmatrix},\qquad \alpha_2=\begin{bmatrix}a\\1\end{bmatrix},\qquad x=\begin{bmatrix}x_1\\x_2\end{bmatrix},}
若二次型 $f(x_1,x_2)=\sum\limits_{i=1}^{2}(\alpha_i,x)^2$ 正定，其中 $(\alpha_i,x)$ 表示向量 $\alpha_i$，$x$ 的内积，则 $a$ 的取值范围是 \underline{\hspace{2.5em}}.

% OPD 蓝色标注（原稿）：→$D_{23}$（化归经典形式），下方蓝色提示语「目的：凑二次型定义」、「按 $f=x^{\mathrm T}Ax$ 的方向走」（方法论提示语，已删）

【解】应填 $a\ne\dfrac12$.

\fmll{f(x_1,x_2)=(\alpha_1,x)^2+(\alpha_2,x)^2}
\fmll{\phantom{f(x_1,x_2)}=(x,\alpha_1)(\alpha_1,x)+(x,\alpha_2)(\alpha_2,x)\qquad(\text{利用内积定义}(\alpha,\beta)=(\beta,\alpha))}
\fmll{\phantom{f(x_1,x_2)}=x^{\mathrm T}\alpha_1\alpha_1^{\mathrm T}x+x^{\mathrm T}\alpha_2\alpha_2^{\mathrm T}x=x^{\mathrm T}(\alpha_1\alpha_1^{\mathrm T}+\alpha_2\alpha_2^{\mathrm T})x}
\fmll{\phantom{f(x_1,x_2)}=x^{\mathrm T}\left(\begin{bmatrix}1\\2\end{bmatrix}[1,2]+\begin{bmatrix}a\\1\end{bmatrix}[a,1]\right)x}
\fmll{\phantom{f(x_1,x_2)}=x^{\mathrm T}\left(\begin{bmatrix}1&2\\2&4\end{bmatrix}+\begin{bmatrix}a^2&a\\a&1\end{bmatrix}\right)x}
\fmll{\phantom{f(x_1,x_2)}=x^{\mathrm T}\begin{bmatrix}1+a^2&2+a\\2+a&5\end{bmatrix}x.}

由题意知，$f$ 正定，故 $1+a^2>0$，$\begin{vmatrix}1+a^2&2+a\\2+a&5\end{vmatrix}>0$，即 $(2a-1)^2>0$，于是当且仅当 $a\ne\dfrac12$ 时，$f$ 正定.
\end{fcard}

% -----------------------------------------------------------
% PDF p106（印刷页 95）
% -----------------------------------------------------------

% OPD 蓝色标注（原稿）：（$O_2$（盯住目标 2）$+D_{23}$（化归经典形式））（\fsec 标题中的 OPD 括号已删，保留方法名）

\fsec{二、配方法与正交变换法的异同}

% OPD 蓝色标注（原稿）：页右缘竖排模块名「形式化归体系块」（OPD 体系块标记，已删）

1. 命题语言

（1）配方法.

二次型语言：将 $f=x^{\mathrm T}Ax$ 通过配方法化为标准形，并求出可逆变换矩阵 $C$.

% OPD 蓝色标注（原稿）：→$D_{22}$（转换等价表述）

矩阵语言：求可逆矩阵 $C$，使得 $C^{\mathrm T}AC=\Lambda$. $\quad(\Lambda\ \text{与}\ A\ \text{合同})$

% OPD 蓝色标注（原稿）：蓝色提示语「要掌握这两个等价的语言表达并能相互转化」（方法论提示语，已删）

\begin{fnote}
【注】考研数学中的配方法指“拉格朗日配方法”，约定简称为配方法.
\end{fnote}

（2）正交变换法.

二次型语言：将 $f=x^{\mathrm T}Ax$ 通过正交变换法化为标准形，并求出正交矩阵 $Q$.

% OPD 蓝色标注（原稿）：→$D_{22}$（转换等价表述）

矩阵语言：求正交矩阵 $Q$，使得 $Q^{\mathrm T}AQ=\Lambda$. $\quad(\Lambda\ \text{与}\ A\ \text{既合同又相似})$

% OPD 蓝色标注（原稿）：蓝色提示语「要掌握这两个等价的语言表达并能相互转化」（方法论提示语，已删）

2. 过程与结果的异同

% OPD 蓝色标注（原稿）：版心左侧蓝色提示语「两种方法一定要搞清，不要只会套公式」（方法论提示语，已删）

$A^{\mathrm T}=A$ 条件下，

\fmll{f(x)=x^{\mathrm T}Ax=}

①配方法（可逆线性变换）：$x=Cy$，$C$ 可逆.

使得 $f\overset{x=Cy}{=}y^{\mathrm T}\Lambda y$，其中 $C^{\mathrm T}AC=\Lambda$（使 $A$ 合同于对角矩阵）.

$C$ 不一定由 $A$ 的特征向量拼成，$\Lambda$ 也不一定是由 $A$ 的特征值拼成的.

②正交变换法（可逆线性变换）：$x=Qy$（这里的 $Q$ 不仅可逆，还满足 $Q^{-1}=Q^{\mathrm T}$），使得 $f\overset{x=Qy}{=}y^{\mathrm T}\Lambda y$，其中 $Q^{\mathrm T}AQ=Q^{-1}AQ=\Lambda$.

二者区别：在配方法中，$C$ 只满足可逆，所以 $C^{-1}$ 不一定等于 $C^{\mathrm T}$，但是在正交变换法中，变换手段 $Q$ 满足 $Q^{-1}=Q^{\mathrm T}$.

二者相同点：它们的正、负惯性指数是对应相等的.

\begin{fcard}{例9.3}
已知矩阵 $A=\begin{bmatrix}4&1&-2\\1&1&1\\-2&1&a\end{bmatrix}$ 与 $B=\begin{bmatrix}k&0&0\\0&6&0\\0&0&0\end{bmatrix}$ 合同.

（1）求 $a$ 的值及 $k$ 的取值范围；

（2）若存在正交矩阵 $Q$ 使得 $Q^{\mathrm T}AQ=B$，求 $k$ 及 $Q$.

% OPD 蓝色标注（原稿）：→$D_{22}$（转换等价表述），准确写出定义

【解】（1）由已知，存在可逆矩阵 $P$ 使得 $P^{\mathrm T}BP=A$，又 $|B|=0$，故 $|A|=3a-12=0$，解得 $a=4$.

$A$ 对应的二次型为

\fmll{f(x_1,x_2,x_3)=4x_1^2+x_2^2+4x_3^2+2x_1x_2-4x_1x_3+2x_2x_3}
\fmll{\phantom{f(x_1,x_2,x_3)}=\left(2x_1+\frac12x_2-x_3\right)^2+\frac34(x_2+2x_3)^2,}

% OPD 蓝色标注（原稿）：蓝色提示语「$B$ 的特征值为 $k,6,0$，$|B|=k\cdot6\cdot0=0$；$A,B$ 合同，故 $A$ 有一个特征值 $0$」（数学说明，已按数学旁注排入正文）

% ---- PDF p107（印刷页 96）----

故经可逆线性变换 $\begin{cases}y_1=2x_1+\frac12x_2-x_3,\\[2pt]y_2=x_2+2x_3,\\[2pt]y_3=x_3,\end{cases}$ 得 $f=y_1^2+\frac34y_2^2$.

$\quad(1\ \text{和}\ \frac34\ \text{不一定是特征值})$

故 $f$ 的正惯性指数为 2. 因为 $A$ 与 $B$ 合同，所以 $B$ 对应二次型的正惯性指数为 2. 因此，$k$ 的取值范围是 $(0,+\infty)$.

（2）因为

\fml{|\lambda E-A|=\begin{vmatrix}\lambda-4&-1&2\\-1&\lambda-1&-1\\2&-1&\lambda-4\end{vmatrix}=\lambda(\lambda-3)(\lambda-6),}

所以 $A$ 的特征值为 $\lambda_1=3$，$\lambda_2=6$，$\lambda_3=0$.

% OPD 蓝色标注（原稿）：→$D_{22}$（转换等价表述），由 $Q^{\mathrm T}=Q^{-1}$，得相似的定义

因为存在正交矩阵 $Q$ 使得 $Q^{\mathrm T}AQ=B$，所以 $A$ 与 $B$ 相似，它们有相同的特征值，于是 $k=3$.

当 $\lambda_1=3$ 时，解方程组 $(3E-A)x=0$，得单位特征向量 $\eta_1=\dfrac{1}{\sqrt3}[1,1,1]^{\mathrm T}$；

当 $\lambda_2=6$ 时，解方程组 $(6E-A)x=0$，得单位特征向量 $\eta_2=\dfrac{1}{\sqrt2}[-1,0,1]^{\mathrm T}$；

当 $\lambda_3=0$ 时，解方程组 $(0E-A)x=0$，得单位特征向量 $\eta_3=\dfrac{1}{\sqrt6}[1,-2,1]^{\mathrm T}$.

% OPD 蓝色标注（原稿）：蓝色旁注「实对称矩阵不同特征值对应的特征向量正交，故只需单位化」（与上方正文说明重复，留痕即可）

令

\fml{Q=[\eta_1,\eta_2,\eta_3]=\begin{bmatrix}\dfrac{1}{\sqrt3}&-\dfrac{1}{\sqrt2}&\dfrac{1}{\sqrt6}\\[8pt]\dfrac{1}{\sqrt3}&0&-\dfrac{2}{\sqrt6}\\[8pt]\dfrac{1}{\sqrt3}&\dfrac{1}{\sqrt2}&\dfrac{1}{\sqrt6}\end{bmatrix},}

则 $Q$ 是正交矩阵，且 $Q^{\mathrm T}AQ=\begin{bmatrix}3&0&0\\0&6&0\\0&0&0\end{bmatrix}=B$.
\end{fcard}

3. 惯性指数

\begin{fcard}{例9.4}
$f(x_1,x_2,x_3)=-2x_1x_2-2x_1x_3+6x_2x_3$ 的正惯性指数为（$\quad$）.

（A）3\qquad（B）2\qquad（C）1\qquad（D）0

【解】应选（C）.

令 $\begin{cases}x_1=y_1+y_2,\\[2pt]x_2=y_1-y_2,\\[2pt]x_3=y_3,\end{cases}$ 则

% OPD 蓝色标注（原稿）：$D_{23}$（化归经典形式），蓝色提示语「没有平方项，创造平方项」（方法论提示语，已删）

\fmll{f=-2y_1^2+2y_2^2+4y_1y_3-8y_2y_3}
\fmll{\phantom{f}=-2(y_1-y_3)^2+2(y_2-2y_3)^2-6y_3^2,}
\end{fcard}

% -----------------------------------------------------------
% PDF p108（印刷页 97）
% -----------------------------------------------------------

再令 $\begin{cases}z_1=y_1-y_3,\\[2pt]z_2=y_2-2y_3,\\[2pt]z_3=y_3,\end{cases}$ 则

\fml{f=-2z_1^2+2z_2^2-6z_3^2,}

故 $f$ 的正惯性指数为 1.

\begin{fcard}{例9.5}
（仅数学一）$f(x_1,x_2,x_3)=x_1x_2+x_1x_3-3x_2x_3=1$ 表示（$\quad$）.

（A）椭球面\qquad（B）双曲柱面

（C）双叶双曲面\qquad（D）单叶双曲面

【解】应选（D）.

法一\quad 令 $\begin{cases}y_1=x_2+x_3,\\[2pt]y_2=x_2-x_3,\\[2pt]y_3=x_1,\end{cases}$ 则

\fmll{f(x_1,x_2,x_3)=y_3y_1-3\dfrac{y_1+y_2}{2}\times\dfrac{y_1-y_2}{2}}
\fmll{\phantom{f(x_1,x_2,x_3)}=-\dfrac34\left(y_1^2-y_2^2\right)+y_1y_3}
\fmll{\phantom{f(x_1,x_2,x_3)}=\dfrac34y_2^2-\dfrac34\left(y_1^2-\dfrac43y_1y_3\right)}
\fmll{\phantom{f(x_1,x_2,x_3)}=\dfrac34y_2^2-\dfrac34\left(y_1-\dfrac23y_3\right)^2+\dfrac13y_3^2=1,}

令 $\begin{cases}z_1=\dfrac{\sqrt3}{2}y_2,\\[8pt]z_2=\dfrac{1}{\sqrt3}y_3,\\[8pt]z_3=\dfrac{\sqrt3}{2}\left(y_1-\dfrac23y_3\right),\end{cases}$ 则 $f(x_1,x_2,x_3)=z_1^2+z_2^2-z_3^2=1$，且作的两次变换均为可逆线性变换.

故 $f(x_1,x_2,x_3)=x_1x_2+x_1x_3-3x_2x_3=1$ 表示的是单叶双曲面.

法二\quad 二次型 $f(x_1,x_2,x_3)=x_1x_2+x_1x_3-3x_2x_3$ 的矩阵

\fml{A=\begin{bmatrix}0&\dfrac12&\dfrac12\\[8pt]\dfrac12&0&-\dfrac32\\[8pt]\dfrac12&-\dfrac32&0\end{bmatrix},}

% ---- PDF p109（印刷页 98）----

由

\fml{|\lambda E-A|=\begin{vmatrix}\lambda&-\dfrac12&-\dfrac12\\[8pt]-\dfrac12&\lambda&\dfrac32\\[8pt]-\dfrac12&\dfrac32&\lambda\end{vmatrix}=0.}

得 $A$ 的特征值为 $\lambda_1=\dfrac32$，$\lambda_2=\dfrac{-3+\sqrt{17}}{4}$，$\lambda_3=\dfrac{-3-\sqrt{17}}{4}$，即 $\lambda_1>0$，$\lambda_2>0$，$\lambda_3<0$，故 $f(x_1,x_2,x_3)=\dfrac32y_1^2+\dfrac{-3+\sqrt{17}}{4}y_2^2+\dfrac{-3-\sqrt{17}}{4}y_3^2=1$ 表示的是单叶双曲面.
\end{fcard}

% OPD 蓝色标注（原稿）：（$O_3$（盯住目标 3）$+D_{23}$（化归经典形式））（\fsec 标题中的 OPD 括号已删，保留方法名）
% OPD 蓝色标注（原稿）：版心左缘竖排模块名「形式化归体系块」（OPD 体系块标记，已删）

\fsec{三、伪配方法}

“平方和式 $A^2+B^2+C^2$” 未必就是（拉格朗日）配方法得来的结果，故若非拉格朗日配方法，则称伪配方法. 要注意伪配方法的变换矩阵是否可逆.

①如果变换没有可逆性，则有可能改变表达式的几何性质，如封闭性，此时不能得出平方和式正定；

②如果变换是可逆的，则平方和式正定.

\begin{fcard}{例9.6}
设二次型 $f(x_1,x_2,x_3)=(x_1+x_2)^2+(x_2+x_3)^2+(ax_1+x_3)^2$ 正定，则 $a$ 的取值范围是 \underline{\hspace{2.5em}}.

【解】应填 $a\ne-1$.

由于 $f(x_1,x_2,x_3)\ge 0$，且 $f(x_1,x_2,x_3)=0\iff\begin{cases}x_1+x_2=0,\\[2pt]x_2+x_3=0,\\[2pt]ax_1+x_3=0,\end{cases}$ 则方程组的系数行列式为

\fml{(-1)\ \text{倍加至}\quad\begin{vmatrix}1&1&0\\0&1&1\\a&0&1\end{vmatrix}=\begin{vmatrix}1&1&0\\0&1&1\\0&-1&a\end{vmatrix}=a+1.}

①当 $a+1\ne0$ 时，$f(x_1,x_2,x_3)=0\iff x_1=x_2=x_3=0$，故当 $a\ne-1$ 时，$f(x_1,x_2,x_3)$ 正定.（$f>0\iff\begin{bmatrix}x_1\\x_2\\x_3\end{bmatrix}\ne 0$）

②当 $a=-1$ 时，存在 $\begin{bmatrix}x_1\\x_2\\x_3\end{bmatrix}\ne 0$，使得 $f$ 等于 0（与正定的定义相悖）.

% OPD 蓝色标注（原稿）：本页右缘蓝色行外注（承接下式）「当 $a=-1$ 时」（数学推导，保留）

当 $a=-1$ 时，

$f=(x_1+x_2)^2+(x_2+x_3)^2+(x_1-x_3)^2$，

令 $\begin{cases}x_1+x_2=y_1,\\[2pt]x_2+x_3=y_2,\\[2pt]x_1-x_3=y_3,\end{cases}$ 即 $\begin{bmatrix}1&1&0\\0&1&1\\1&0&-1\end{bmatrix}\begin{bmatrix}x_1\\x_2\\x_3\end{bmatrix}=\begin{bmatrix}y_1\\y_2\\y_3\end{bmatrix}$

因为其行列式为 $0$，所以此矩阵不可逆，故此变换为不可逆线性变换，会改变原变量间的性质，也即原图形为非封闭的图形，但是通过这种不可逆的线性变换图形变成封闭的了，也是把整个问题的性质改变了
\end{fcard}

% ---- PDF p110（印刷页 99）----

% 原稿此处的青色【注】框自 PDF p109 末页起、跨页至 PDF p110 首，为独立于例 9.6 的注释框

\begin{fnote}
【注】对于 $f(x_1,x_2,x_3)=(a_1x_1+a_2x_2+a_3x_3)^2+(b_1x_1+b_2x_2+b_3x_3)^2+(c_1x_1+c_2x_2+c_3x_3)^2$ 的情形，可总结如下做题方法：

令 $f=0$，即 $\begin{cases}a_1x_1+a_2x_2+a_3x_3=0,\\[2pt]b_1x_1+b_2x_2+b_3x_3=0,\\[2pt]c_1x_1+c_2x_2+c_3x_3=0,\end{cases}$ 计算 $|A|=\begin{vmatrix}a_1&a_2&a_3\\b_1&b_2&b_3\\c_1&c_2&c_3\end{vmatrix}$，若 $|A|\ne 0$，则 $f$ 正定；若 $|A|=0$，则 $f$ 不正定.
\end{fnote}

% OPD 蓝色标注（原稿）：（$O_4$（盯住目标 4）$+D_1$（常规操作）$+D_{23}$（化归经典形式））（\fsec 标题中的 OPD 括号已删，保留方法名）

\fsec{四、正交变换法的传递性}

若 $A$ 相似于 $B$，$B$ 相似于 $C$，则 $A$ 相似于 $C$. 这里 $B$ 常为 $\Lambda$.

\begin{fcard}{例9.7}
二次型 $f(x_1,x_2,x_3)=x_1^2-x_2x_3$ 经正交变换 $x=Qy$ 化为二次型

\fml{g(y_1,y_2,y_3)=y_1y_2+ay_3^2.}

（1）求 $a$ 的值；

（2）求正交矩阵 $Q$.

% OPD 蓝色标注（原稿）：例 9.7 题干右侧「→基本局面」（方法论旁注，已删）

【解】（1）二次型 $f(x_1,x_2,x_3)$ 与 $g(y_1,y_2,y_3)$ 的矩阵分别为

\fml{A=\begin{bmatrix}1&0&0\\[4pt]0&0&-\dfrac12\\[8pt]0&-\dfrac12&0\end{bmatrix},\qquad B=\begin{bmatrix}0&\dfrac12&0\\[8pt]\dfrac12&0&0\\[4pt]0&0&a\end{bmatrix},}

且 $Q^{\mathrm T}AQ=B$. 因为 $Q$ 为正交矩阵，所以矩阵 $A$ 与 $B$ 相似，故 $\operatorname{tr}(A)=\operatorname{tr}(B)$，从而 $a=1$.

（2）由于 $|\lambda E-A|=(\lambda-1)\left(\lambda-\dfrac12\right)\left(\lambda+\dfrac12\right)$，因此矩阵 $A$ 与 $B$ 的特征值均为 1，$\dfrac12$，$-\dfrac12$，易求得矩阵 $A$ 对应于特征值 1，$\dfrac12$，$-\dfrac12$ 的特征向量依次为 $\begin{bmatrix}1\\0\\0\end{bmatrix}$，$\begin{bmatrix}0\\-1\\1\end{bmatrix}$，$\begin{bmatrix}0\\1\\1\end{bmatrix}$，又因为实对称矩阵属于不同特征值的特征向量相互正交，故单位化即可，依次为 $\begin{bmatrix}1\\0\\0\end{bmatrix}$，$\begin{bmatrix}0\\-\dfrac{1}{\sqrt2}\\[6pt]\dfrac{1}{\sqrt2}\end{bmatrix}$，$\begin{bmatrix}0\\\dfrac{1}{\sqrt2}\\[6pt]\dfrac{1}{\sqrt2}\end{bmatrix}$

令 $Q_1=\begin{bmatrix}1&0&0\\[4pt]0&-\dfrac{1}{\sqrt2}&\dfrac{1}{\sqrt2}\\[8pt]0&\dfrac{1}{\sqrt2}&\dfrac{1}{\sqrt2}\end{bmatrix}$，则 $Q_1$ 为正交矩阵，且 $Q_1^{\mathrm T}AQ_1=\begin{bmatrix}1&0&0\\[4pt]0&\dfrac12&0\\[4pt]0&0&-\dfrac12\end{bmatrix}$.

% ---- PDF p111（印刷页 100）----

同理，可求得矩阵 $B$ 对应于特征值 1，$\dfrac12$，$-\dfrac12$ 的单位特征向量依次为 $\begin{bmatrix}0\\0\\1\end{bmatrix}$，$\begin{bmatrix}\dfrac{1}{\sqrt2}\\[6pt]\dfrac{1}{\sqrt2}\\[6pt]0\end{bmatrix}$，$\begin{bmatrix}-\dfrac{1}{\sqrt2}\\[6pt]\dfrac{1}{\sqrt2}\\[6pt]0\end{bmatrix}$，

令 $Q_2=\begin{bmatrix}0&\dfrac{1}{\sqrt2}&-\dfrac{1}{\sqrt2}\\[8pt]0&\dfrac{1}{\sqrt2}&\dfrac{1}{\sqrt2}\\[8pt]1&0&0\end{bmatrix}$，则 $Q_2$ 为正交矩阵，且

\fml{Q_2^{\mathrm T}BQ_2=\begin{bmatrix}1&0&0\\[4pt]0&\dfrac12&0\\[4pt]0&0&-\dfrac12\end{bmatrix}.}

% OPD 蓝色标注（原稿）：$D_{23}$（化归经典形式）

由于 $Q_1^{\mathrm T}AQ_1=Q_2^{\mathrm T}BQ_2$，因此 $Q_2Q_1^{\mathrm T}AQ_1Q_2^{\mathrm T}=B$，从而 $Q=Q_1Q_2^{\mathrm T}=\begin{bmatrix}0&0&1\\-1&0&0\\0&1&0\end{bmatrix}$ 为所求正交矩阵.
\end{fcard}

% OPD 蓝色标注（原稿）：（$O_5$（盯住目标 5）$+D_1$（常规操作）$+D_{23}$（化归经典形式））（\fsec 标题中的 OPD 括号已删，保留方法名）

\fsec{深化二　化标准形的三种方法}

\begin{fdeep}{正交变换 vs 配方法：结论不同，别用错}
把二次型化为标准形有两条路，\textbf{它们给出的标准形一般不同}：

\begin{center}
\begin{tabular}{@{}lll@{}}
\toprule
方法 & 变换类型 & 标准形系数 \\
\midrule
正交变换 $x=Qy$ & 正交矩阵（$Q^{\mathrm T}Q=E$） & \textbf{恰为 $A$ 的全部特征值} \\
配方法 $x=Cy$ & 一般可逆矩阵 & 只保证\textbf{正负惯性指数相同}，系数不是特征值 \\
\bottomrule
\end{tabular}
\end{center}

\textbf{关键区别}：正交变换是\textbf{合同且相似}，所以系数必须是特征值；
配方法\textbf{只是合同}，系数可以任意（只要正负个数对）。
因此：
\begin{itemize}[leftmargin=2.2em,itemsep=0.22em]
\item 题目问“\textbf{求 }$A$\textbf{ 的特征值}”，或要求“正交变换”→ \textbf{必须用正交变换法}，
      系数就是特征值，算错说明特征值算错了；
\item 题目只说“\textbf{化二次型为标准形}”（未指明方法）→ \textbf{两种都可以}，
      此时用配方法通常更快（不必求特征向量、不必正交化）；
\item 题目问“\textbf{求正负惯性指数}”→ 配方法足够（本来就是不变量）。
\end{itemize}
\textbf{一个有用的自检}：若用正交变换算出的标准形系数之和 $\ne\operatorname{tr}A$，
说明特征值算错了（因为 $\sum\lambda_i=\operatorname{tr}A$，第 7 章）。
这个自检几乎不花时间，却能拦住大部分计算失误。

\textbf{规范形是唯一的}：无论用哪种方法，继续把系数化为 $\pm1$
（正系数变 $+1$、负系数变 $-1$），得到的\textbf{规范形唯一}——这正是惯性定理的内容。

\fbref{}服务考纲〔二次型〕“掌握用正交变换化二次型为标准形的方法，会用配方法化二次型为标准形”。
考纲要求\textbf{两种方法都会}，而它们的适用场合不同：
\textbf{要特征值就必须用正交变换}（配方法给不出特征值），
\textbf{只要标准形/惯性指数就优先配方法}（更快）。
混用会导致“算出漂亮的系数却不是特征值”这类隐性错误，
用 $\sum\text{系数}=\operatorname{tr}A$ 自检即可当场发现。
\end{fdeep}

\begin{fdeep}{配方法的两种情形与操作步骤}
配方法的本质是把一个变量一次配完、再对剩下的变量重复（对变量个数作归纳）。
\textbf{情形一：有平方项，不妨 $a_{11}\ne0$。} 把含 $x_1$ 的项全部集中配成完全平方：
\fml{f=a_{11}\Bigl(x_1+\sum_{j=2}^{n}a_{11}^{-1}a_{1j}x_j\Bigr)^{2}
+\sum_{i=2}^{n}\sum_{j=2}^{n}b_{ij}x_ix_j}
后半部分只含 $x_2,\dots,x_n$，是 $n-1$ 元二次型。取替换
$y_1=x_1+\sum_{j\ge2}a_{11}^{-1}a_{1j}x_j,\ y_2=x_2,\ \dots,\ y_n=x_n$ 即可。
该替换的系数矩阵是单位上三角阵、行列式为 $1$，故一定\textbf{非退化}。
\textbf{情形二：没有平方项（所有 $a_{ii}=0$）但有交叉项 $a_{ij}\ne0\ (i\ne j)$。}
先“制造”一个平方项，不妨设 $a_{12}\ne0$，取
\fml{x_1=z_1+z_2,\qquad x_2=z_1-z_2,\qquad x_k=z_k\ \ (k=3,\dots,n)}
则 $f=2a_{12}z_1^{2}-2a_{12}z_2^{2}+\dots$，其中 $z_1^{2}$ 的系数 $2a_{12}\ne0$，化为情形一。
（连交叉项都没有时，$f$ 本身已是 $n-1$ 元二次型，直接归纳。）
\textbf{考场动作}：情形一不必固定配 $x_1$，优先选“有平方项且系数最简单”的那个变量；
情形二先作上面的和差替换，把交叉项拆成平方项再回到情形一。

\fbref{}：服务考纲〔二次型〕“会用配方法化二次型为标准形”。配方是\textbf{唯一不必求特征值、不必正交化}的方法，纯代数操作，速度最快。学生的唯一卡点就是“二次型不含平方项”——记住 $x_1=z_1+z_2,\ x_2=z_1-z_2$ 这一步（其行列式为 $-2\ne0$，必非退化），任何二次型都能拉回情形一。
\end{fdeep}

\begin{fdeep}{初等变换法：一行一列成对变换，同时记录变换矩阵}
除配方法与正交变换法之外的第三条路：对 $n$ 阶对称矩阵 $A$ 施行一系列\textbf{成对}的初等变换
——作一次列变换就立刻作一次同样的行变换——这正好是合同变换 $P^{\mathrm T}AP$。
把它写成上下两块的分块矩阵，全程只需盯住上块：
\fml{\begin{pmatrix}A\\E\end{pmatrix}\ \longrightarrow\ \begin{pmatrix}\Lambda\\C\end{pmatrix}}
（其中 $\Lambda$ 为对角阵，$C$ 为可逆阵）
操作规则：上块 $A$ 每作一次\textbf{列}变换，下块 $E$ 就跟着作一次\textbf{同样的列}变换；
上块再补作一次对应的\textbf{行}变换（下块不动）。做完后 $C^{\mathrm T}AC=\Lambda$，
令 $x=Cy$ 即得标准形 $f=y^{\mathrm T}\Lambda y=\lambda_1y_1^{2}+\dots+\lambda_ny_n^{2}$。
\textbf{与配方法的差别}：配方法只交出替换、不管矩阵；初等变换法把\textbf{可逆矩阵 $C$ 一并算出来}，
所以题目一旦要求“写出所作的非退化线性替换”，它比配方更稳、更不易出错。

\fbref{}：服务考纲〔二次型〕“会用配方法化二次型为标准形”“掌握用正交变换化二次型为标准形”。当题干要求\textbf{同时给出标准形与所用替换 $x=Cy$}（$C$ 只要求可逆、不必正交）时，初等变换法是标准解法：下块忠实记录了每一步，$C$ 是“白送”的。唯一要守住的一条纪律——\textbf{下块永远只作列变换}，写成行变换就前功尽弃。
\end{fdeep}

\begin{fdeep}{顺序主子式法：不求特征值也能写出标准形}
设 $n$ 阶实对称矩阵 $A$ 的各阶顺序主子式 $\Delta_1,\Delta_2,\dots,\Delta_n$ \textbf{全不为零}
（约定 $\Delta_0=1$）。则存在\textbf{非退化}线性替换 $X=CY$ 使
\fml{f=\frac{\Delta_1}{\Delta_0}y_1^{2}+\frac{\Delta_2}{\Delta_1}y_2^{2}
+\dots+\frac{\Delta_n}{\Delta_{n-1}}y_n^{2}}
即第 $k$ 个系数恰为 $\lambda_k=\Delta_k/\Delta_{k-1}$，并且还可要求 $C$ 是单位上三角矩阵。
证明思路（对阶数归纳）：把 $A$ 分块为 $\begin{pmatrix}A_{n-1}&\beta\\ \beta^{\mathrm T}&a_{nn}\end{pmatrix}$，
由归纳假设有 $S$ 使 $S^{\mathrm T}A_{n-1}S=\operatorname{diag}(\lambda_1,\dots,\lambda_{n-1})$；
再用一个单位上三角的合同变换消去右上角 $\beta$，即得对角阵
$\operatorname{diag}(\lambda_1,\dots,\lambda_{n-1},\lambda_n)$，两边取行列式就有
$\lambda_n=\Delta_n/\Delta_{n-1}$。
\textbf{使用前提}：所有 $\Delta_k\ne0$；若某个 $\Delta_k=0$，本法失效，须先作一次合同置换再试。

\fbref{}：服务考纲〔二次型〕“掌握用正交变换化二次型为标准形，会用配方法化二次型为标准形”。这是\textbf{第三种、也是最省事}的求标准形办法：只算 $n$ 个顺序主子式，不用求特征值、不用求特征向量、不用配方；而 $\Delta_k/\Delta_{k-1}$ 的符号直接给出正负惯性指数（看相邻两个 $\Delta$ 是否同号），判正定、判合同一步到位。
\end{fdeep}

\fsec{五、合同的判定与手段}

1. 同阶实对称矩阵 $A$，$B$ 合同的判定

用正、负惯性指数：$A$，$B$ 合同 $\iff p_A=p_B$，$q_A=q_B$（相同的正、负惯性指数）.

2. 已知 $A$，$\Lambda$（$\Lambda$ 是对角矩阵），求可逆矩阵 $C$，使得 $C^{\mathrm T}AC=\Lambda$

\begin{fcard}{例9.8}
已知 $A=\begin{bmatrix}3&2&1\\2&2&1\\1&1&1\end{bmatrix}$，$\Lambda=\begin{bmatrix}2&0&0\\0&3&0\\0&0&1\end{bmatrix}$，求可逆矩阵 $C$，使得 $C^{\mathrm T}AC=\Lambda$.

【解】①配方.

\fmll{f=x^{\mathrm T}Ax=3x_1^2+2x_2^2+x_3^2+4x_1x_2+2x_1x_3+2x_2x_3}
\fmll{\phantom{f}=(x_1^2+2x_1x_2+x_2^2+2x_1x_3+2x_2x_3+x_3^2)+(x_1^2+2x_1x_2+x_2^2)+x_1^2}
\fmll{\phantom{f}=x_1^2+(x_1+x_2)^2+(x_1+x_2+x_3)^2}
\fmll{\phantom{f}=2\cdot\left(\frac{x_1}{\sqrt2}\right)^2+3\cdot\left(\frac{x_1+x_2}{\sqrt3}\right)^2+1\cdot(x_1+x_2+x_3)^2.}

% OPD 蓝色标注（原稿）：版心左缘蓝色提示语「配方时，不一定要先配 $x_1$，再配 $x_2$，看哪个最容易配，就先配哪个」（数学方法说明，已按下行内联排入正文）
% OPD 蓝色标注（原稿）：蓝色提示语「盯着 $A$ 的对角线元素，平方项分别提出 $2,3,1$」（数学方法说明，已按下行内联排入正文）

$\quad(\text{配方时，不一定要先配}\ x_1\ \text{，再配}\ x_2\ \text{，看哪个最容易配，就先配哪个})$

$\quad(\text{盯着}\ A\ \text{的对角线元素，平方项分别提出}\ 2,3,1)$

% -----------------------------------------------------------
% PDF p112（印刷页 101）
% -----------------------------------------------------------

②换元.

令 $\begin{cases}y_1=\dfrac{x_1}{\sqrt2},\\[8pt]y_2=\dfrac{x_1+x_2}{\sqrt3},\\[8pt]y_3=x_1+x_2+x_3,\end{cases}$ 于是有 $\begin{bmatrix}y_1\\y_2\\y_3\end{bmatrix}=\begin{bmatrix}\dfrac{1}{\sqrt2}&0&0\\[8pt]\dfrac{1}{\sqrt3}&\dfrac{1}{\sqrt3}&0\\[8pt]1&1&1\end{bmatrix}\begin{bmatrix}x_1\\x_2\\x_3\end{bmatrix}$.

③求逆.

记 $x=Cy$，其中 $C=\begin{bmatrix}\dfrac{1}{\sqrt2}&0&0\\[8pt]\dfrac{1}{\sqrt3}&\dfrac{1}{\sqrt3}&0\\[8pt]1&1&1\end{bmatrix}^{-1}=\begin{bmatrix}\sqrt2&0&0\\-\sqrt2&\sqrt3&0\\0&-\sqrt3&1\end{bmatrix}$，则 $f=x^{\mathrm T}Ax=(Cy)^{\mathrm T}A(Cy)=y^{\mathrm T}C^{\mathrm T}ACy=y^{\mathrm T}\Lambda y$，

即矩阵 $C$ 可使 $C^{\mathrm T}AC=\Lambda$.
\end{fcard}

3. 已知 $A$，$B$（$B$ 不是对角矩阵），求可逆矩阵 $C$，使得 $C^{\mathrm T}AC=B$

\begin{fcard}{例9.9}
已知实矩阵 $A=\begin{bmatrix}2&2\\2&1\end{bmatrix}$，$B=\begin{bmatrix}4&3\\3&1\end{bmatrix}$，且 $A$ 与 $B$ 合同. 求可逆矩阵 $D$，使得 $A=D^{\mathrm T}BD$.

% OPD 蓝色标注（原稿）：例 9.9 题干右侧「→变换手段」（方法论旁注，已删）

【解】①对 $f$ 配方、换元，写 $D_1$.

对于

\fml{f(x_1,x_2)=x^{\mathrm T}Ax=2(x_1+x_2)^2-x_2^2,}

令 $\begin{cases}z_1=x_1+x_2,\\[2pt]z_2=x_2,\end{cases}$ 即 $\begin{bmatrix}z_1\\z_2\end{bmatrix}=\begin{bmatrix}1&1\\0&1\end{bmatrix}\begin{bmatrix}x_1\\x_2\end{bmatrix}=D_1\begin{bmatrix}x_1\\x_2\end{bmatrix}$，则 $f(x_1,x_2)=2z_1^2-z_2^2$.

②对 $g$ 配方、换元，写 $D_2$.

对于

\fml{g(y_1,y_2)=y^{\mathrm T}By=4\left(y_1+\dfrac34y_2\right)^2-\dfrac54y_2^2,}

令 $\begin{cases}z_1=\sqrt2y_1+\dfrac{3\sqrt2}{4}y_2,\\[8pt]z_2=\dfrac{\sqrt5}{2}y_2,\end{cases}$ 即 $\begin{bmatrix}z_1\\z_2\end{bmatrix}=\begin{bmatrix}\sqrt2&\dfrac{3\sqrt2}{4}\\[8pt]0&\dfrac{\sqrt5}{2}\end{bmatrix}\begin{bmatrix}y_1\\y_2\end{bmatrix}=D_2\begin{bmatrix}y_1\\y_2\end{bmatrix}$，则 $g(y_1,y_2)=2z_1^2-z_2^2$.

③令 $D_1x=D_2y$，求 $D_2^{-1}D_1$.

于是有 $D_1\begin{bmatrix}x_1\\x_2\end{bmatrix}=D_2\begin{bmatrix}y_1\\y_2\end{bmatrix}$，故 $\begin{bmatrix}y_1\\y_2\end{bmatrix}=D_2^{-1}D_1\begin{bmatrix}x_1\\x_2\end{bmatrix}=D\begin{bmatrix}x_1\\x_2\end{bmatrix}$，即

\fml{D=\begin{bmatrix}\sqrt2&\dfrac{3\sqrt2}{4}\\[8pt]0&\dfrac{\sqrt5}{2}\end{bmatrix}^{-1}\begin{bmatrix}1&1\\0&1\end{bmatrix}=\begin{bmatrix}\dfrac{1}{\sqrt2}&-\dfrac{3\sqrt5}{10}\\[8pt]0&\dfrac{2}{\sqrt5}\end{bmatrix}\begin{bmatrix}1&1\\0&1\end{bmatrix}=\begin{bmatrix}\dfrac{1}{\sqrt2}&\dfrac{1}{\sqrt2}-\dfrac{3\sqrt5}{10}\\[8pt]0&\dfrac{2}{\sqrt5}\end{bmatrix},}

% 例 9.9 的收尾结论（「故 $A=D^{\mathrm T}BD$」）在 PDF p113（印刷页 102），由后续片段承接。
\end{fcard}

\fsec{深化三　惯性定理、规范形与惯性指数}

\begin{fdeep}{惯性定理的证明思路：为什么正负惯性指数是不变量}
同一个实二次型可以化成许多不同的标准形，但规范形只有一个：
\fml{f=d_1y_1^{2}+\dots+d_ry_r^{2}\ (d_i\ne0)\ \Longrightarrow\
f=z_1^{2}+\dots+z_p^{2}-z_{p+1}^{2}-\dots-z_r^{2}}
（实数域中正数总可开平方，故可把每个 $d_i$ 变成 $\pm1$）。关键只在证明 $p$ 唯一。
\textbf{反证法}：设同一 $f$ 经 $X=BY$ 与 $X=CZ$ 分别化成规范形，要证 $p=q$。若 $p>q$，
令 $G=C^{-1}B=(g_{ij})$，则 $Z=GY$。作齐次线性方程组（前 $q$ 个方程取自 $Z=GY$）：
\fml{\begin{cases}
g_{11}y_1+g_{12}y_2+\dots+g_{1n}y_n=0,\\
\dots\dots\dots\dots\dots\dots\\
g_{q1}y_1+g_{q2}y_2+\dots+g_{qn}y_n=0,\\
y_{p+1}=0,\ \dots,\ y_n=0.
\end{cases}}
它含 $n$ 个未知量、而方程只有 $q+(n-p)=n-(p-q)<n$ 个，故有非零解 $(k_1,\dots,k_n)$。
把该解代入规范形的左端，得 $k_1^{2}+\dots+k_p^{2}>0$；代入右端，因它是方程组前 $q$ 个方程的解
故 $z_1=\dots=z_q=0$，得 $-z_{q+1}^{2}-\dots-z_r^{2}\le0$。矛盾。所以 $p\le q$，同理 $q\le p$，
于是 $p=q$，规范形唯一。

\fbref{}：服务考纲〔二次型〕“了解…标准形规范形概念及惯性定理”。“方程个数少于未知量个数 $\Longrightarrow$ 有非零解”是惯性定理的\textbf{唯一考法}：凡是要证“两个由二次型符号定义的数相等”，都从齐次方程组入手；它也顺带说明了正负惯性指数为什么是\textbf{合同不变量}（只依赖二次型本身，与替换无关）。
\end{fdeep}

\begin{fdeep}{规范形唯一、标准形不唯一：复域与实域两套规范形}
同一个二次型用不同的非退化线性替换，得到的标准形系数可以完全不同：例如某个三元二次型
既可化为 $2y_1^{2}-2y_2^{2}+6y_3^{2}$，也可化为 $2w_1^{2}-\frac{1}{2}w_2^{2}+\frac{2}{3}w_3^{2}$。
所以\textbf{标准形不唯一}，能唯一确定下来的只是“非零系数的个数”和“正负系数的个数”。
把非零系数继续用 $y_i=\frac{1}{\sqrt{d_i}}z_i$ 化成 $\pm1$，得到两套\textbf{唯一}的规范形：
复二次型的规范形为
\fml{z_1^{2}+z_2^{2}+\dots+z_r^{2}}
实二次型的规范形为
\fml{z_1^{2}+\dots+z_p^{2}-z_{p+1}^{2}-\dots-z_r^{2}}
（依据是复数总能开平方、正实数总能开平方。）由此得到两条\textbf{合同判据}：
复对称矩阵 $A$ 与 $B$ 合同 $\iff$ 秩相等；
实对称矩阵 $A$ 与 $B$ 合同 $\iff$ 正、负惯性指数分别相等（等价于秩相同且正惯性指数相同）。

\fbref{}：服务考纲〔二次型〕“了解二次型秩、合同变换与合同矩阵、标准形规范形概念及惯性定理”。考场上问“两个二次型的标准形是否相同”，答“标准形相同”是\textbf{错}的：标准形不唯一；只有\textbf{规范形}才唯一，而“规范形相同 $\iff$ 秩相同且正惯性指数相同”。这一句把“标准形 / 规范形”这对极易混淆的概念一次性分清。
\end{fdeep}

\begin{fdeep}{惯性指数就是特征值的符号计数}
对实对称矩阵 $A$（$r=\operatorname{rank}A$）：\textbf{正惯性指数 $p$ 等于正特征值的个数}，
\textbf{负惯性指数 $r-p$ 等于负特征值的个数}，秩等于非零特征值的个数。理由是正交变换下
\fml{Q^{\mathrm T}AQ=Q^{-1}AQ=\operatorname{diag}(\lambda_1,\dots,\lambda_n)}
它\textbf{既是合同又是相似}：合同保惯性指数、相似保特征值，两边的 $\pm1$ 计数必须对上。
由此得几条好用的推论：
（1）$A$ 正定 $\iff$ 全部特征值 $>0$；半正定 $\iff$ 全部特征值 $\ge0$（且至少有一个为 $0$）；
（2）符号差 $=2p-r=\sum_{i=1}^{n}\operatorname{sgn}\lambda_i$；
（3）正负惯性指数可“数特征值符号”，也可“看 $\Delta_k/\Delta_{k-1}$ 的符号”（即相邻两个顺序主子式
是否同号），\textbf{两条路算出的答案必须一致}。

\fbref{}：服务考纲〔二次型〕“掌握二次型及其矩阵表示”“掌握用正交变换化二次型为标准形”，并与〔矩阵的特征值和特征向量〕打通。求正负惯性指数有三条路：数正特征值、看顺序主子式是否变号、配方法数正系数。\textbf{用两条不同的路算同一个 $p$}，是二次型部分最便宜的错误拦截手段（若用正交变换算出的标准形系数之和与 $\operatorname{tr}A$ 不符，也立刻暴露特征值算错）。
\end{fdeep}

\begin{fdeep}{惯性指数的估计：先把二次型写成 $p$ 个平方减 $q$ 个平方}
设 $f=l_1^{2}+\dots+l_p^{2}-l_{p+1}^{2}-\dots-l_{p+q}^{2}$，其中每个 $l_i$ 都是
$x_1,\dots,x_n$ 的\textbf{一次齐次式}。则 $f$ 的\textbf{正惯性指数 $\le p$，负惯性指数 $\le q$}。
证明思路（反证 + 齐次方程组，与惯性定理同一套路）：设 $f$ 经 $X=TY$ 化为标准形
$y_1^{2}+\dots+y_r^{2}-y_{r+1}^{2}-\dots$，若 $r>p$，则方程组
（前 $p$ 个取 $l_i=0$，后 $n-r$ 个取 $y_j=0$，共 $p+(n-r)<n$ 个方程）必有非零解 $X_0$。
此时左端 $f(X_0)=-l_{p+1}^{2}-\dots-l_{p+q}^{2}\le0$，右端却是 $y_1^{2}+\dots+y_r^{2}>0$，矛盾。
\textbf{用途}：数一数已给的平方项个数，就能立刻给出正惯性指数的\textbf{上界}。

\fbref{}：服务考纲〔二次型〕“理解正定二次型、正定矩阵的概念，并掌握其判别法”。见到 $f=\sum(\ \cdot\ )^{2}-\sum(\ \cdot\ )^{2}$ 这种结构（含参数题尤其常见），先用本结论卡住 $p$ 的上界，往往一步就排除一半选项；它也是判断“$f$ 是否可能正定”最快的反驳工具（平方项不足 $n$ 个即不正定）。
\end{fdeep}

\begin{fdeep}{二次型能分解为两个一次型之积的充要条件}
非零实二次型 $f(x_1,\dots,x_n)$ 可以写成
\fml{f=(u_1x_1+\dots+u_nx_n)(v_1x_1+\dots+v_nx_n)}
的\textbf{充要条件}是：$f$ 的\textbf{秩为 $1$}，或者\textbf{秩为 $2$ 且符号差为 $0$}。
换成惯性指数的说法：正负惯性指数只能是 $(1,0)$、$(0,1)$（秩 $1$，即 $f$ 是正/负定的一次型平方）
或 $(1,1)$（秩 $2$ 且符号差 $0$）。证明思路：若 $u,v$ 线性相关，不妨 $u=kv$，
作替换 $y_1=v_1x_1+\dots+v_nx_n,\ y_i=x_i\ (i\ge2)$ 得 $f=ky_1^{2}$，秩为 $1$；
若 $u,v$ 线性无关，令 $y_1=\sum u_ix_i,\ y_2=\sum v_ix_i$ 得 $f=y_1y_2$，
再令 $y_1=z_1+z_2,\ y_2=z_1-z_2$ 得 $f=z_1^{2}-z_2^{2}$，正是秩 $2$、符号差 $0$。
反之把 $y$ 用 $C^{-1}$ 反解回 $x$ 即可构造出两个一次型。

\fbref{}：服务考纲〔二次型〕“了解二次型秩…及惯性定理”。“$f$ 能否分解为两个一次因式之积”是反复出现的判断题，硬展开系数矩阵很慢；\textbf{只要求出秩与符号差，答案立刻出来}：秩 $1$ 行；秩 $2$ 看符号差是否为 $0$；秩 $\ge3$ 一定不行。这是惯性指数最“解题化”的一个用处。
\end{fdeep}

\begin{fdeep}{两类高频特殊二次型的秩与符号差}
（1）$f=ayz+bzx+cxy$（$a,b,c$ 为实数）：
$a,b,c$ 全为零时，秩 $=0$、符号差 $=0$；$abc=0$ 但不全为零时，秩 $=2$、符号差 $=0$；
$abc>0$ 时，秩 $=3$、符号差 $=-1$；$abc<0$ 时，秩 $=3$、符号差 $=1$。
（推导：$c\ne0$ 时 $A$ 合同于 $\operatorname{diag}\bigl(c,-c,-\frac{ab}{c}\bigr)$，再按 $c$ 与 $ab$ 的符号分类。）
（2）$f=n\sum_{i=1}^{n}x_i^{2}-\Bigl(\sum_{i=1}^{n}x_i\Bigr)^{2}$，它还可以写成离差平方和
\fml{f=\sum_{1\le i<j\le n}(x_i-x_j)^{2}}
所以它是\textbf{半正定}二次型，秩 $=$ 正惯性指数 $=n-1$，负惯性指数为 $0$，
规范形是 $y_1^{2}+\dots+y_{n-1}^{2}$。（它的矩阵 $A=nE-\alpha\alpha^{\mathrm T}$，
$\alpha=(1,1,\dots,1)^{\mathrm T}$，特征值为 $n$（$n-1$ 重）与 $0$。）

\fbref{}：服务考纲〔二次型〕“了解二次型秩…及定义与惯性定理”。这两类结构在填空、选择里反复出现：“$axy+byz+czx$ 的秩与符号差”按 $abc$ 的符号\textbf{直接背出四种答案}；“$n\sum x_i^{2}-(\sum x_i)^{2}$”要能一眼看出它是离差平方和，从而半正定、秩 $n-1$、负惯性指数 $0$，完全不必展开求特征值。
\end{fdeep}

% ===========================================================
% 《线性代数9讲》第 9 讲  二次型 —— 片段 B（本讲后半）
% 源：线代9讲强化.pdf  PDF p113–p121（印刷页 102–110），共 9 页
%
% 处理说明：
%   1) OPD 方法论标记按 AGENTS.md §7.1 决定二「不保留方法论」的统一口径剔除：
%      · \fsec 标题里的 OPD 括号连同其中的 OPD 词组**整段删除**：
%        「六、合同与相似的异同」（原附 $O_6$（盯住目标 6）+$D_1$（常规操作）
%          +$D_{23}$（化归经典形式））——已删；
%        「七、正定的判定与应用」（原附 $O_7$（盯住目标 7）+$D_1$（常规操作）
%          +$D_{23}$（化归经典形式））——已删；
%        「八、二次型的最大值」（原附 $O_8$（盯住目标 8）+$D_1$（常规操作）
%          +$D_{22}$（转换等价表述）+$D_{23}$（化归经典形式））——已删；
%      · 正文蓝色纯 OPD 代码 / 方法论语句一律以 `% OPD 蓝色标注（原稿）：…` 留痕，不排入正文；
%      · 版心左侧竖排模块名「等价表述体系块」、右侧竖排模块名「隐含条件体系块」
%        （PDF p115，OPD 体系块标记）→ 删除并留痕；
%      · 原稿各页右上角蓝色装饰性二维码（PDF p113、p115、p116、p118，非正文）→ 未转写。
%   2) **数学操作旁注保留**（原稿蓝色手写，照原文以内联 `\quad(\text{…})` 排入）：
%      · PDF p114（103）$f=z_1^2+z_2^2$ 右侧「$1$ 不一定是特征值」；
%      · PDF p116（105）例 9.13 题设右侧「已知 $D^{\mathrm T}D$（格拉姆矩阵），求 $D$」；
%      · PDF p118（107）八节推导框首行上方「$Q$ 为正交矩阵」；
%      · PDF p119（108）（2）式计算末「转化成瑞利商形式，故只需分析矩阵 $M$ 的最大特征值即可」。
%   3) 页眉（「第9讲 二次型」）、页脚（102–110 及云纹）为每页重复装饰，按规范忽略。
%   4) 本片 9 页**无「习题」栏**（已逐页核实 PDF p113–p121）；PDF p121 为全书最后一页，**整页空白**。
%   5) 交底：PDF p120（109）的【注】框是**印刷正文**（非 \fsec 标题、非蓝色旁注），
%      依「严禁增删内容」照抄；其中含「三向解题法」及 $D_{23}$（化归经典形式）等 OPD 方法论措辞，
%      已在最终回复中报告，供主控决定是否再删。
%
% 接缝说明：
%   · 本片首行「则 $A=D^{\mathrm T}BD$.」承 PDF p112（印刷页 101）上一例题【解】的末行；
%   · 例 9.11 的解答跨 PDF p113～p114（102～103），例 9.15 的解答跨 PDF p118～p119（107～108），
%     均未另开卡片，以正文承接。
% ===========================================================

% -----------------------------------------------------------
% PDF p113（印刷页 102）
% -----------------------------------------------------------
% 承 PDF p112（印刷页 101）上一例题【解】的末行
则 $A=D^{\mathrm{T}}BD$.

\begin{fnote}
【注】此处得到 $C^{\mathrm{T}}AC=B$，即 $A$ 与 $B$ 合同，$C$ 可逆，但 $A$ 与 $B$ 未必相似，故不一定存在正交矩阵 $Q$，使得 $Q^{\mathrm{T}}AQ=B$. 例 9.11 会再次研究这个问题.
\end{fnote}

\fsec{六、合同与相似的异同}
% OPD 蓝色标注（原稿）：节标题下附（$O_6$（盯住目标 6）+$D_1$（常规操作）+$D_{23}$（化归经典形式）），按统一口径整段删除

对于实对称矩阵 $A$ 与 $B$，相似必合同，反之不成立.

\begin{fcard}{例9.10}
设矩阵 $A=\begin{bmatrix}2&-1&-1\\-1&2&-1\\-1&-1&2\end{bmatrix}$，$B=\begin{bmatrix}1&0&0\\0&1&0\\0&0&0\end{bmatrix}$，则 $A$ 与 $B$（\quad）.

（A）合同且相似\qquad（B）合同，但不相似

（C）不合同，但相似\qquad（D）既不合同，也不相似

【解】应选（B）.

因为
\fml{|\lambda E-A|=\begin{vmatrix}\lambda-2&1&1\\1&\lambda-2&1\\1&1&\lambda-2\end{vmatrix}=\lambda(\lambda-3)^2,}
所以矩阵 $A$ 的特征值为 $3$，$3$，$0$，由此可知矩阵 $A$ 与 $B$ 不相似，从而选项（A）和（C）错误.

又因为实对称矩阵 $A$ 相似且合同于对角矩阵
\fml{C=\begin{bmatrix}3&0&0\\0&3&0\\0&0&0\end{bmatrix},}
而矩阵 $C$ 显然合同于矩阵 $B$，根据合同关系的传递性知矩阵 $A$ 合同于 $B$，即选项（B）正确.
\end{fcard}

\begin{fcard}{★★★例9.11}
已知二次型
\fml{f(x_1,x_2,x_3)=x_1^2+2x_2^2+2x_3^2+2x_1x_2-2x_1x_3,}
\fml{g(y_1,y_2,y_3)=y_1^2+y_2^2+y_3^2+2y_2y_3,}
（1）求可逆变换 $x=Py$，将 $f(x_1,x_2,x_3)$ 化成 $g(y_1,y_2,y_3)$.

（2）是否存在正交变换 $x=Qy$，将 $f(x_1,x_2,x_3)$ 化成 $g(y_1,y_2,y_3)$？

【解】（1）由于
\fml{f(x_1,x_2,x_3)=x_1^2+2x_2^2+2x_3^2+2x_1x_2-2x_1x_3=(x_1+x_2-x_3)^2+(x_2+x_3)^2,}
于是作可逆线性变换

% -----------------------------------------------------------
% PDF p114（印刷页 103）—— 例 9.11 续；（本页末为七节起始）
% -----------------------------------------------------------
\fml{\begin{bmatrix}z_1\\z_2\\z_3\end{bmatrix}=\begin{bmatrix}1&1&-1\\0&1&1\\0&0&1\end{bmatrix}\begin{bmatrix}x_1\\x_2\\x_3\end{bmatrix},}
得 $f(x_1,x_2,x_3)=z_1^2+z_2^2$\quad（$1$ 不一定是特征值）.
% 数学操作旁注（原稿蓝色手写）：上式右侧箭头 →「$1$ 不一定是特征值」，已按规范内联保留

由于 $g(y_1,y_2,y_3)=y_1^2+y_2^2+y_3^2+2y_2y_3=y_1^2+(y_2+y_3)^2$，于是作可逆线性变换
\fml{\begin{bmatrix}z_1\\z_2\\z_3\end{bmatrix}=\begin{bmatrix}1&0&0\\0&1&1\\0&0&1\end{bmatrix}\begin{bmatrix}y_1\\y_2\\y_3\end{bmatrix},}
得 $g(y_1,y_2,y_3)=z_1^2+z_2^2$.

令
\fml{P=\begin{bmatrix}1&1&-1\\0&1&1\\0&0&1\end{bmatrix}^{-1}\begin{bmatrix}1&0&0\\0&1&1\\0&0&1\end{bmatrix}=\begin{bmatrix}1&-1&1\\0&1&0\\0&0&1\end{bmatrix},}
则在可逆变换 $x=Py$ 下，
\fml{f(x_1,x_2,x_3)=g(y_1,y_2,y_3).}

（2）二次型 $f(x_1,x_2,x_3)$ 与 $g(y_1,y_2,y_3)$ 对应的矩阵分别为
\fml{A=\begin{bmatrix}1&1&-1\\1&2&0\\-1&0&2\end{bmatrix},\qquad B=\begin{bmatrix}1&0&0\\0&1&1\\0&1&1\end{bmatrix}.}
由于 $\operatorname{tr}(A)\ne\operatorname{tr}(B)$，因此矩阵 $A$ 与 $B$ 不相似，故不存在正交变换 $x=Qy$，将 $f(x_1,x_2,x_3)$ 化成 $g(y_1,y_2,y_3)$.
\end{fcard}

\begin{fnote}
【注】以上两题，就是从命题手法上改变了同一知识点的难度，命制成例 9.10 这样的选择题，题干要求明确、直白，考生一般答得很好，但命制成例 9.11 这样的解答题，第（1）问是给出“$f$ 与 $g$ 合同”的等价说法；第（2）问是给出“$f$ 与 $g$ 所对应的二次型矩阵 $A$ 与 $B$ 不相似”的等价说法，且问题呈开放性，给考生答题带来了不小的难度.
\end{fnote}

\fsec{深化四　正定性与半正定性}

\begin{fdeep}{正定的六条等价判别：按题目给的条件挑最省事的一条}
设 $A$ 为 $n$ 阶实对称矩阵，以下命题\textbf{两两等价}：
\begin{enumerate}[leftmargin=2.2em,itemsep=0.22em]
\item \key{定义}：对任意 $x\ne0$ 都有 $x^{\mathrm T}Ax>0$；
\item \textbf{顺序主子式}全为正：$\lvert A_k\rvert>0\ (k=1,2,\dots,n)$，$A_k$ 为左上角 $k$ 阶子阵；
\item \textbf{特征值}全为正：$\lambda_i>0$；
\item \textbf{正惯性指数} $p=n$（即规范形为 $y_1^{2}+\cdots+y_n^{2}$）；
\item $A$ 与单位矩阵\textbf{合同}：存在可逆 $C$ 使 $A=C^{\mathrm T}EC=C^{\mathrm T}C$；
\item 存在\textbf{可逆}矩阵 $C$ 使 $A=C^{\mathrm T}C$。
\end{enumerate}

\textbf{怎么选（考场动作）}：题目给什么就用哪条——
\begin{itemize}[leftmargin=2.2em,itemsep=0.22em]
\item 给的是\textbf{具体数字矩阵} → 用 ②（算顺序主子式，最快，只需算 $n$ 个行列式）；
\item 给的是\textbf{特征值}或需要求特征值 → 用 ③；
\item 给的是\textbf{抽象条件}（如 $A=B^{\mathrm T}B$、$A$ 可逆）→ 用 ⑥ 或 ⑤；
\item 证明题要“说明某个二次型正定” → 从 ① 出发（配方或直接估计）最稳。
\end{itemize}

\textbf{两条必须区分的“否命题”}：
\begin{itemize}[leftmargin=2.2em,itemsep=0.22em]
\item \textbf{“主对角元全正”不能推出正定}。反例：$A=\begin{pmatrix}1&2\\2&1\end{pmatrix}$
      对角元全为 $1>0$，但特征值为 $3$ 与 $-1$，\textbf{不正定}。
      正确方向只有“正定 $\Rightarrow$ 主对角元全正”（取 $x=e_i$ 即得 $a_{ii}>0$）。
\item \textbf{“所有主子式全正”才是正定}（不只是顺序主子式）；
      而“顺序主子式全正”与“所有主子式全正”对\textbf{对称矩阵}是等价的——
      这就是为什么判正定只需算顺序主子式。
\end{itemize}

\warn{半正定的相应判别（这里有个必须注意的差别）}：
\begin{itemize}[leftmargin=2.2em,itemsep=0.22em]
\item \textbf{正定}：只需查\textbf{顺序}主子式（因为对称性使“所有主子式全正”与“顺序主子式全正”等价）；
\item \textbf{半正定}：必须查\textbf{所有}主子式 $\ge0$（含非顺序的），
      只查顺序主子式\textbf{会漏判}。
\end{itemize}
\warn{反例（务必记住）}：$A=\begin{pmatrix}1&1&0\\1&1&0\\0&0&-1\end{pmatrix}$。
它的顺序主子式为 $1,\ 0,\ 0$，\textbf{全部 $\ge0$}；但特征值为 $-1,\ 0,\ 2$，含负值，
故\textbf{不是半正定}（它连正定都不是）。原因在于“所有主子式”中还含有负的那个
（取第 3 行第 3 列的单元素主子式 $\lvert A_{33}\rvert=-1<0$）。
\must{所以判半正定的口诀是：“所有主子式非负”，而不是“顺序主子式非负”。}

\fbref{}服务考纲〔二次型〕“理解正定二次型、正定矩阵的概念，并掌握其判别法”。
六条要\must{整体记住}，因为选择题常问“下列哪个是 $A$ 正定的充要条件”；
而“按题目条件挑一条”能显著省时间（数字矩阵用主子式、抽象条件用 $C^{\mathrm T}C$）。
两个否命题（对角元正 $\ne$ 正定、顺序主子式要配合对称性）是最常见的陷阱。
\end{fdeep}

\begin{fdeep}{负定、半正定、半负定、不定的判别}
设 $f(x)=x^{\mathrm T}Ax$（$A$ 实对称，$x\ne0$）：$f>0$ 称\textbf{正定}；$f<0$ 称\textbf{负定}；
$f\ge0$ 称\textbf{半正定}；$f\le0$ 称\textbf{半负定}；既有 $f(x_1)>0$ 又有 $f(x_2)<0$ 称\textbf{不定}。
\textbf{负定没有独立的判别法}——加个负号化成正定即可：$f$ 负定 $\iff$ $-f$ 正定 $\iff$ $-A$ 正定，
于是它的顺序主子式满足
\fml{(-1)^{k}\Delta_k>0\qquad(k=1,2,\dots,n)}
即\textbf{奇数阶顺序主子式为负、偶数阶为正}。
\textbf{不定}的判定最省事：只要能举出一个 $x$ 使 $f(x)>0$、另一个使 $f(x)<0$ 即可；
或看特征值有正有负、正负惯性指数都不为 $0$。此外还有一条好用的等价说法：
$A$ 半正定或半负定 $\iff$ 凡满足 $\alpha^{\mathrm T}A\alpha=0$ 的 $\alpha$ 都有 $A\alpha=0$
（必要条件显然；充分性用反证：若正负惯性指数都 $\ge1$，取规范形下 $y_{p+1}=1$、其余为 $0$，
得 $\alpha^{\mathrm T}A\alpha=0$ 而 $A\alpha\ne0$）。

\fbref{}：服务考纲〔二次型〕“理解正定二次型、正定矩阵的概念，并掌握其判别法”。考纲字面只写“正定”，但选择题常考负定与不定。两条口诀：\textbf{负定 = 加负号变正定，故判据为 $(-1)^{k}\Delta_k>0$}；\textbf{判不定只需举两个反例向量}。最后那条“$\alpha^{\mathrm T}A\alpha=0\Longrightarrow A\alpha=0$”是判半正定/半负定的巧法。
\end{fdeep}

\begin{fdeep}{半正定的判别：为什么“顺序主子式非负”不够}
$n$ 元实二次型 $f=x^{\mathrm T}Ax$（$A$ 实对称）半正定，有以下\textbf{等价条件}：
（1）$f$ 半正定；（2）正惯性指数 $p$ 等于秩 $r$（即负惯性指数为 $0$）；
（3）存在可逆 $C$ 使 $C^{\mathrm T}AC=\operatorname{diag}(d_1,\dots,d_n)$，其中 $d_i\ge0$；
（4）存在实矩阵 $C$ 使 $A=C^{\mathrm T}C$；
（5）$A$ 的\textbf{所有主子式}（行指标与列指标相同）都 $\ge0$。
\textbf{关键差别}：正定只需查\textbf{顺序}主子式，半正定必须查\textbf{所有}主子式。
\textbf{反例}：$f(x_1,x_2)=-x_2^{2}$，其矩阵为
$\begin{pmatrix}0&0\\0&-1\end{pmatrix}$，顺序主子式为 $0,\ 0$ 全部 $\ge0$，
但 $f(0,1)=-1<0$，\textbf{不是半正定}——漏掉的是一阶主子式 $a_{22}=-1$。
（证明思路：必要性取 $x$ 只在前 $k$ 个分量上变动；充分性反证，若正负惯性指数都 $\ge1$，
取 $C^{\mathrm T}AC=\operatorname{diag}(E_p,-E_q,O)$ 下 $y_{p+1}=1$、其余为 $0$ 的那个向量。）

\fbref{}：服务考纲〔二次型〕“掌握正定矩阵的判别法”（半正定是其自然延伸）。这是二次型部分\warn{最经典的陷阱}：题干给“顺序主子式都 $\ge0$”，一定要警惕它只对正定充分。一句话记牢：\textbf{正定看顺序主子式，半正定看全部主子式}；判半正定还有个更快的写法——把它写成 $C^{\mathrm T}C$ 或若干平方和。
\end{fdeep}

\begin{fdeep}{正定矩阵的元素、行列式与合同结论}
$A$ 正定（实对称）时，以下都是可以\textbf{直接引用}的结论：
（1）$|A|>0$，所有对角元 $a_{ii}>0$；更强的有
\fml{|A|\le a_{11}a_{22}\cdots a_{nn},\qquad |A|\le a_{nn}\Delta_{n-1}}
（$\Delta_{n-1}$ 为 $n-1$ 阶顺序主子式）；
（2）$A$ 中\textbf{绝对值最大的元素必在主对角线上}：由二阶主子式 $a_{ii}a_{jj}-a_{ij}^{2}>0$
得 $|a_{ij}|<\sqrt{a_{ii}a_{jj}}\le a_{kk}$，其中 $a_{kk}=\max\{a_{11},\dots,a_{nn}\}$；
（3）$A$ 与 $A^{-1}$ \textbf{合同}（因 $A^{-1}=(A^{-1})^{\mathrm T}AA^{-1}$），
所以 $x^{\mathrm T}Ax$ 与 $x^{\mathrm T}A^{-1}x$ 有相同的正负惯性指数；
（4）对 $m\times n$ 实矩阵 $B$（$m\ge n$），$B^{\mathrm T}AB$ 正定 $\iff$ $\operatorname{rank}B=n$；
（5）$A$ 可逆 $\iff$ 存在矩阵 $B$ 使 $AB+B^{\mathrm T}A$ 正定。

\fbref{}：服务考纲〔二次型〕“理解正定二次型、正定矩阵的概念，并掌握其判别法”，并与〔矩阵〕的逆与秩打通。(2) 把“找最大元”这类选择题一步解决；(3)(4) 是“抽象二次型判正定”的固定套路——见 $B^{\mathrm T}AB$ 就想 $\operatorname{rank}B=n$；(1) 是行列式估计题的入口。
\end{fdeep}

\begin{fdeep}{同时合同对角化：$A$ 正定时可以把 $B$ 一起化为对角阵}
设 $A$ 是 $n$ 阶\textbf{正定}矩阵，$B$ 是 $n$ 阶\textbf{实对称}矩阵，则存在可逆矩阵 $P$ 使
\fml{P^{\mathrm T}AP=E,\qquad P^{\mathrm T}BP=\operatorname{diag}(\mu_1,\dots,\mu_n)}
其中 $\mu_1,\dots,\mu_n$ 恰是方程 $|\lambda A-B|=0$ 的 $n$ 个根，且全为实数。
证法：$A$ 正定 $\Longrightarrow$ 存在可逆 $C$ 使 $C^{\mathrm T}AC=E$；此时 $C^{\mathrm T}BC$ 仍实对称，
故存在\textbf{正交}矩阵 $Q$ 使 $Q^{\mathrm T}(C^{\mathrm T}BC)Q=\operatorname{diag}(\mu_i)$，
取 $P=CQ$（正交变换不会破坏已化好的 $E$）。两个常用推论：
（1）$A$ 正定、$B$ 实对称时，存在可逆变换 $x=Py$ 把 $x^{\mathrm T}Ax$ 化为\textbf{规范形}、
同时把 $x^{\mathrm T}Bx$ 化为\textbf{标准形}；
（2）若 $|A-\lambda B|=0$ 有 $n$ 个\textbf{互异}实根，则存在非退化 $X=CY$ 同时把两个二次型化为标准形。

\fbref{}：服务考纲〔二次型〕“掌握用正交变换化二次型为标准形的方法”。“两型同时化简”是二次型大题的常见收尾要求。固定套路是\textbf{先把正定的那个 $A$ 化为 $E$}（用配方或初等变换都行），再对剩下的实对称阵作正交对角化；顺序反了就做不出来，$\mu_i$ 也一定要交代清楚是 $|\lambda A-B|=0$ 的根。
\end{fdeep}

\fsec{七、正定的判定与应用}
% OPD 蓝色标注（原稿）：节标题下附（$O_7$（盯住目标 7）+$D_1$（常规操作）+$D_{23}$（化归经典形式）），按统一口径整段删除

\begin{fcard}{1. 前提}
\fml{A=A^{\mathrm{T}}\text{（}A\text{ 是实对称矩阵）}.}
\end{fcard}

% -----------------------------------------------------------
% PDF p115（印刷页 104）
% -----------------------------------------------------------
% OPD 蓝色标注（原稿）：→$D_2$（脱胎换骨）（箭头指向下条标题「2. 二次型 $f=x^{\mathrm T}Ax$ 正定的充要条件」）
% OPD 蓝色标注（原稿）：版心左侧竖排模块名「等价表述体系块」（方法论模块标签，已删）

\begin{fcard}{2. 二次型 $f=x^{\mathrm{T}}Ax$ 正定的充要条件}
$n$ 元二次型 $f=x^{\mathrm{T}}Ax$ 正定
\fml{\left\{\begin{aligned}
&\iff \text{对任意的 } x\ne0\text{，有 } x^{\mathrm{T}}Ax>0\text{（定义）}\\
&\iff A\text{ 的特征值 }\lambda_i>0\ (i=1,2,\dots,n)\\
&\iff f\text{ 的正惯性指数 }p=n\\
&\iff \text{存在可逆矩阵 } D\text{，使得 } A=D^{\mathrm{T}}D\\
&\iff A\text{ 与 } E\text{ 合同}\\
&\iff A\text{ 的各阶顺序主子式均大于 }0.
\end{aligned}\right.}
% 数学操作旁注（原稿蓝色手写）：上式第 4 个「$\iff$」右侧箭头 →「尤其注意」（指向「存在可逆矩阵 $D$，使得 $A=D^{\mathrm T}D$」）
\end{fcard}

\begin{fcard}{3. 二次型 $f=x^{\mathrm{T}}Ax$ 正定的必要条件}
①$a_{ii}>0\ (i=1,2,\dots,n)$.

②$|A|>0$.
\end{fcard}

% OPD 蓝色标注（原稿）：版心右侧竖排模块名「隐含条件体系块」（方法论模块标签，已删；起自「4. 重要结论」行）

\begin{fcard}{4. 重要结论}
①若 $A$ 正定，则 $A^{-1}$，$A^{*}$，$A^{m}$（$m$ 为正整数），$kA$（$k>0$），$C^{\mathrm{T}}AC$（$C$ 可逆）均正定.

②若 $A$，$B$ 正定，则 $A+B$ 正定，$\begin{bmatrix}A&O\\O&B\end{bmatrix}$ 正定.

③若 $A$，$B$ 正定，则 $AB$ 正定的充要条件是 $AB=BA$.
\end{fcard}

\begin{fcard}{例9.12}
设矩阵
\fml{A=\begin{bmatrix}1&2\\-2&-a\end{bmatrix},\qquad B=\begin{bmatrix}1&0\\1&a\end{bmatrix},}
若 $f(x,y)=|xA+yB|$ 是正定二次型，则 $a$ 的取值范围是（\quad）.

（A）$(0,2-\sqrt{3})$\qquad（B）$(2-\sqrt{3},2+\sqrt{3})$

（C）$(2+\sqrt{3},4)$\qquad（D）$(0,4)$

% OPD 蓝色标注（原稿）：→$D_1$（常规操作），表达式“新”并不意味着内容“新”，其一下看看便知

【解】应选（B）.

\fml{xA+yB=\begin{bmatrix}x&2x\\-2x&-ax\end{bmatrix}+\begin{bmatrix}y&0\\y&ay\end{bmatrix}=\begin{bmatrix}x+y&2x\\-2x+y&-ax+ay\end{bmatrix},}
\fml{\begin{aligned}
f(x,y)=|xA+yB|&=\begin{vmatrix}x+y&2x\\-2x+y&-ax+ay\end{vmatrix}=-a(x+y)(x-y)-2x(-2x+y)\\
&=(4-a)x^2-2xy+ay^2,
\end{aligned}}
二次型的矩阵为 $\begin{bmatrix}4-a&-1\\-1&a\end{bmatrix}$，若 $f(x,y)=|xA+yB|$ 是正定二次型，则

% -----------------------------------------------------------
% PDF p116（印刷页 105）—— 例 9.12 收尾；例 9.13
% -----------------------------------------------------------
\fml{4-a>0,\qquad \begin{vmatrix}4-a&-1\\-1&a\end{vmatrix}=4a-a^2-1>0,}
解得 $2-\sqrt{3}<a<2+\sqrt{3}$. 选（B）.
\end{fcard}

\begin{fcard}{例9.13}
若可逆矩阵 $D$ 满足
\fml{D^{\mathrm{T}}D=\begin{bmatrix}1&-1&1\\-1&2&-3\\1&-3&6\end{bmatrix},}
则 $D=$ \underline{\hspace{4em}}\quad（已知 $D^{\mathrm{T}}D$（格拉姆矩阵），求 $D$）.
% 数学操作旁注（原稿蓝色手写）：上式右侧箭头 →「已知 $D^{\mathrm T}D$（格拉姆矩阵），求 $D$」，已按规范内联保留

【解】应填 $\begin{bmatrix}1&-1&1\\0&1&-2\\0&0&1\end{bmatrix}$.

法一\quad 记 $\begin{bmatrix}1&-1&1\\-1&2&-3\\1&-3&6\end{bmatrix}=A$，则其对应的二次型为
\fmll{\begin{aligned}
f(x_1,x_2,x_3)&=x^{\mathrm{T}}Ax=x_1^2+2x_2^2+6x_3^2-2x_1x_2+2x_1x_3-6x_2x_3\\
&=(x_1-x_2+x_3)^2+(x_2-2x_3)^2+x_3^2\\
&=[x_1-x_2+x_3,\ x_2-2x_3,\ x_3]\begin{bmatrix}x_1-x_2+x_3\\x_2-2x_3\\x_3\end{bmatrix}\\
&=[x_1,x_2,x_3]\begin{bmatrix}1&0&0\\-1&1&0\\1&-2&1\end{bmatrix}\begin{bmatrix}1&-1&1\\0&1&-2\\0&0&1\end{bmatrix}\begin{bmatrix}x_1\\x_2\\x_3\end{bmatrix}\\
&\overset{\text{记}}{=}(x_1,x_2,x_3)D^{\mathrm{T}}D\begin{bmatrix}x_1\\x_2\\x_3\end{bmatrix},
\end{aligned}}
其中 $A=D^{\mathrm{T}}D$，$D=\begin{bmatrix}1&-1&1\\0&1&-2\\0&0&1\end{bmatrix}$.

法二\quad 在法一中，$f(x_1,x_2,x_3)=x^{\mathrm{T}}Ax=(x_1-x_2+x_3)^2+(x_2-2x_3)^2+x_3^2$，令 $\begin{cases}x_1-x_2+x_3=y_1,\\x_2-2x_3=y_2,\\x_3=y_3,\end{cases}$ 则
$f(x_1,x_2,x_3)=y_1^2+y_2^2+y_3^2$，也即作可逆线性变换 $x=Cy$，其中 $\begin{bmatrix}1&-1&1\\0&1&-2\\0&0&1\end{bmatrix}\begin{bmatrix}x_1\\x_2\\x_3\end{bmatrix}=\begin{bmatrix}y_1\\y_2\\y_3\end{bmatrix}$，得 $C=\begin{bmatrix}1&-1&1\\0&1&-2\\0&0&1\end{bmatrix}^{-1}$.

因此
\fml{f=x^{\mathrm{T}}Ax=y^{\mathrm{T}}C^{\mathrm{T}}ACy=y^{\mathrm{T}}Ey,}

% -----------------------------------------------------------
% PDF p117（印刷页 106）—— 例 9.13 收尾；例 9.14
% -----------------------------------------------------------
故 $C^{\mathrm{T}}AC=E$，则 $A=(C^{-1})^{\mathrm{T}}C^{-1}=D^{\mathrm{T}}D$，故 $D=C^{-1}=\begin{bmatrix}1&-1&1\\0&1&-2\\0&0&1\end{bmatrix}$.
\end{fcard}

\begin{fnote}
【注】所求的矩阵 $D$ 不唯一.
\end{fnote}

\begin{fcard}{例9.14}
设矩阵 $A=\begin{bmatrix}a&1&-1\\1&a&-1\\-1&-1&a\end{bmatrix}$.

（1）求正交矩阵 $P$，使 $P^{\mathrm{T}}AP$ 为对角矩阵；

（2）求正定矩阵 $C$，使 $C^2=(a+3)E-A$，其中 $E$ 为 3 阶单位矩阵.

【解】（1）因为 $|\lambda E-A|=\begin{vmatrix}\lambda-a&-1&1\\-1&\lambda-a&1\\1&1&\lambda-a\end{vmatrix}=(\lambda-a+1)^2(\lambda-a-2)$，所以 $A$ 的特征值为 $\lambda_1=\lambda_2=a-1$，$\lambda_3=a+2$.

当 $\lambda_1=\lambda_2=a-1$ 时，解方程组 $[(a-1)E-A]x=0$，得 $A$ 的线性无关的特征向量 $\xi_1=\begin{bmatrix}-1\\1\\0\end{bmatrix}$，$\xi_2=\begin{bmatrix}1\\0\\1\end{bmatrix}$，进
行施密特正交单位化得 $\eta_1=\begin{bmatrix}-\dfrac{\sqrt{2}}{2}\\[2mm]\dfrac{\sqrt{2}}{2}\\[1mm]0\end{bmatrix}$，$\eta_2=\begin{bmatrix}\dfrac{\sqrt{6}}{6}\\[2mm]\dfrac{\sqrt{6}}{6}\\[2mm]\dfrac{\sqrt{6}}{3}\end{bmatrix}$.

当 $\lambda_3=a+2$ 时，解方程组 $[(a+2)E-A]x=0$，得 $A$ 的特征向量 $\xi_3=\begin{bmatrix}-1\\-1\\1\end{bmatrix}$，单位化得 $\eta_3=\begin{bmatrix}-\dfrac{\sqrt{3}}{3}\\[2mm]-\dfrac{\sqrt{3}}{3}\\[2mm]\dfrac{\sqrt{3}}{3}\end{bmatrix}$.

令 $P=[\eta_1,\eta_2,\eta_3]=\begin{bmatrix}-\dfrac{\sqrt{2}}{2}&\dfrac{\sqrt{6}}{6}&-\dfrac{\sqrt{3}}{3}\\[2mm]\dfrac{\sqrt{2}}{2}&\dfrac{\sqrt{6}}{6}&-\dfrac{\sqrt{3}}{3}\\[2mm]0&\dfrac{\sqrt{6}}{3}&\dfrac{\sqrt{3}}{3}\end{bmatrix}$，则

% -----------------------------------------------------------
% PDF p118（印刷页 107）—— 例 9.14 收尾；八节；例 9.15
% -----------------------------------------------------------
\fml{P^{\mathrm{T}}AP=\begin{bmatrix}a-1&0&0\\0&a-1&0\\0&0&a+2\end{bmatrix},}
故 $P$ 为所求正交矩阵.

（2）由（1）知
\fml{(a+3)E-A=(a+3)E-P\begin{bmatrix}a-1&0&0\\0&a-1&0\\0&0&a+2\end{bmatrix}P^{\mathrm{T}}=P\begin{bmatrix}4&0&0\\0&4&0\\0&0&1\end{bmatrix}P^{\mathrm{T}}.}
% OPD 蓝色标注（原稿）：→$D_{23}$（化归经典形式），对矩阵“开方”是对矩阵“开方”的逆向运算，但总的局面仍是相似对角化

令 $C=P\begin{bmatrix}2&0&0\\0&2&0\\0&0&1\end{bmatrix}P^{\mathrm{T}}$，则 $C^2=(a+3)E-A$. 故所求正定矩阵是
\fml{\begin{aligned}
C&=\begin{bmatrix}-\dfrac{\sqrt{2}}{2}&\dfrac{\sqrt{6}}{6}&-\dfrac{\sqrt{3}}{3}\\[2mm]\dfrac{\sqrt{2}}{2}&\dfrac{\sqrt{6}}{6}&-\dfrac{\sqrt{3}}{3}\\[2mm]0&\dfrac{\sqrt{6}}{3}&\dfrac{\sqrt{3}}{3}\end{bmatrix}\begin{bmatrix}2&0&0\\0&2&0\\0&0&1\end{bmatrix}\begin{bmatrix}-\dfrac{\sqrt{2}}{2}&\dfrac{\sqrt{6}}{6}&-\dfrac{\sqrt{3}}{3}\\[2mm]\dfrac{\sqrt{2}}{2}&\dfrac{\sqrt{6}}{6}&-\dfrac{\sqrt{3}}{3}\\[2mm]0&\dfrac{\sqrt{6}}{3}&\dfrac{\sqrt{3}}{3}\end{bmatrix}^{\mathrm{T}}\\
&=\begin{bmatrix}\dfrac{5}{3}&-\dfrac{1}{3}&\dfrac{1}{3}\\[2mm]-\dfrac{1}{3}&\dfrac{5}{3}&\dfrac{1}{3}\\[2mm]\dfrac{1}{3}&\dfrac{1}{3}&\dfrac{5}{3}\end{bmatrix}.
\end{aligned}}
\end{fcard}

\fsec{八、二次型的最大值}
% OPD 蓝色标注（原稿）：节标题下附（$O_8$（盯住目标 8）+$D_1$（常规操作）+$D_{22}$（转换等价表述）+$D_{23}$（化归经典形式）），按统一口径整段删除

% 原稿：节标题下方为蓝色边框的推导框（无编号），以下按提示框排入
\begin{fnote}
求瑞利商 $\dfrac{x^{\mathrm{T}}Ax}{x^{\mathrm{T}}x}$ 的最值：$f=x^{\mathrm{T}}Ax\overset{x=Qy\ (\text{$Q$ 为正交矩阵})}{=}y^{\mathrm{T}}\Lambda y$
\fml{=\lambda_1y_1^2+\lambda_2y_2^2+\lambda_3y_3^2.}
令 $\lambda_{\max}=\max\{\lambda_1,\lambda_2,\lambda_3\}$；$\lambda_{\min}=\min\{\lambda_1,\lambda_2,\lambda_3\}$.
\fml{\lambda_{\min}\underbrace{(y_1^2+y_2^2+y_3^2)}_{x_1^2+x_2^2+x_3^2}\le f\le\lambda_{\max}\underbrace{(y_1^2+y_2^2+y_3^2)}_{x_1^2+x_2^2+x_3^2}\implies \lambda_{\min}\le\frac{f}{x_1^2+x_2^2+x_3^2}=\frac{x^{\mathrm{T}}Ax}{x^{\mathrm{T}}x}\le\lambda_{\max}}
% 数学操作旁注（原稿）：推导框首行上方蓝色手写「$Q$ 为正交矩阵」及其左向箭头（指向 $x=Qy$），已内联于上式
\end{fnote}

\begin{fcard}{例9.15}
设二次型 $f(x_1,x_2)=x_1^2-4x_1x_2+4x_2^2$，$g(x_1,x_2)$ 的二次型矩阵为
\fml{B=\begin{bmatrix}1&-1\\-1&2\end{bmatrix}.}

（1）是否存在可逆矩阵 $D$，使 $B=D^{\mathrm{T}}D$？若存在，求出矩阵 $D$，若不存在，说明理由.

（2）求 $\max\limits_{x\ne0}\dfrac{f(x)}{g(x)}$，其中 $x=\begin{bmatrix}x_1\\x_2\end{bmatrix}$.

% -----------------------------------------------------------
% PDF p119（印刷页 108）—— 例 9.15 解答
% -----------------------------------------------------------
【解】（1）由于 $B^{\mathrm{T}}=B$，$1>0$，$\begin{vmatrix}1&-1\\-1&2\end{vmatrix}>0$，因此 $B$ 为正定矩阵. 又 $g(x_1,x_2)=(x_1-x_2)^2+x_2^2$，令
$\begin{cases}x_1-x_2=y_1,\\x_2=y_2,\end{cases}$ 于是有 $\begin{cases}x_1=y_1+y_2,\\x_2=y_2,\end{cases}$ 即 $x=\begin{bmatrix}x_1\\x_2\end{bmatrix}=\begin{bmatrix}1&1\\0&1\end{bmatrix}\begin{bmatrix}y_1\\y_2\end{bmatrix}=Cy$，使 $x^{\mathrm{T}}Bx=y^{\mathrm{T}}C^{\mathrm{T}}BCy=y^{\mathrm{T}}y$，即
\fml{C^{\mathrm{T}}BC=E,\qquad B=(C^{-1})^{\mathrm{T}}C^{-1}=\begin{bmatrix}1&-1\\0&1\end{bmatrix}^{\mathrm{T}}\begin{bmatrix}1&-1\\0&1\end{bmatrix}=D^{\mathrm{T}}D.}
故存在可逆矩阵 $D=\begin{bmatrix}1&-1\\0&1\end{bmatrix}$，使 $B=D^{\mathrm{T}}D$.
% OPD 蓝色标注（原稿）：→$D_{23}$（化归经典形式）

（2）设 $f(x)$ 的二次型矩阵为 $A$，则
\fml{\frac{f(x)}{g(x)}=\frac{x^{\mathrm{T}}Ax}{x^{\mathrm{T}}Bx}=\frac{x^{\mathrm{T}}Ax}{x^{\mathrm{T}}D^{\mathrm{T}}Dx}=\frac{x^{\mathrm{T}}Ax}{(Dx)^{\mathrm{T}}Dx}\overset{\text{令}Dx=z}{=}\frac{(D^{-1}z)^{\mathrm{T}}A(D^{-1}z)}{z^{\mathrm{T}}z}=\frac{z^{\mathrm{T}}(D^{-1})^{\mathrm{T}}AD^{-1}z}{z^{\mathrm{T}}z},}
其中 $z=\begin{bmatrix}z_1\\z_2\end{bmatrix}$\quad（转化成瑞利商形式，故只需分析矩阵 $M$ 的最大特征值即可）.
% 数学操作旁注（原稿蓝色手写）：上式右侧「转化成瑞利商形式，故只需分析矩阵 $M$ 的最大特征值即可」，已按规范内联保留

又
\fml{\begin{aligned}
(D^{-1})^{\mathrm{T}}AD^{-1}=C^{\mathrm{T}}AC&=\begin{bmatrix}1&0\\1&1\end{bmatrix}\begin{bmatrix}1&-2\\-2&4\end{bmatrix}\begin{bmatrix}1&1\\0&1\end{bmatrix}=\begin{bmatrix}1&-2\\-1&2\end{bmatrix}\begin{bmatrix}1&1\\0&1\end{bmatrix}\\
&=\begin{bmatrix}1&-1\\-1&1\end{bmatrix}\overset{\text{记}}{=}M,
\end{aligned}}
其对应的二次型为 $h(z_1,z_2)=z_1^2+z_2^2-2z_1z_2$，则由
\fml{|\lambda E-M|=\begin{vmatrix}\lambda-1&1\\1&\lambda-1\end{vmatrix}=(\lambda-1)^2-1=\lambda(\lambda-2)=0,}
得 $\lambda_1=2$，$\xi_1=\dfrac{1}{\sqrt{2}}\begin{bmatrix}-1\\1\end{bmatrix}$；$\lambda_2=0$，$\xi_2=\dfrac{1}{\sqrt{2}}\begin{bmatrix}1\\1\end{bmatrix}$，令 $Q=\dfrac{1}{\sqrt{2}}\begin{bmatrix}-1&1\\1&1\end{bmatrix}$，$z=Qr=Q\begin{bmatrix}r_1\\r_2\end{bmatrix}$，则
\fml{\begin{aligned}
z^{\mathrm{T}}Mz=r^{\mathrm{T}}Q^{\mathrm{T}}MQr&=r^{\mathrm{T}}\begin{bmatrix}2&0\\0&0\end{bmatrix}r=2r_1^2\\
&\le 2(r_1^2+r_2^2)\overset{z^{\mathrm{T}}z=r^{\mathrm{T}}Q^{\mathrm{T}}Qr=r^{\mathrm{T}}r}{=}2(z_1^2+z_2^2),
\end{aligned}}
于是 $z^{\mathrm{T}}Mz\le 2z^{\mathrm{T}}z$，即 $\dfrac{z^{\mathrm{T}}Mz}{z^{\mathrm{T}}z}\le 2$，令 $z_0=Q\begin{bmatrix}1\\0\end{bmatrix}$，得 $\dfrac{\begin{bmatrix}1,0\end{bmatrix}\begin{bmatrix}2&0\\0&0\end{bmatrix}\begin{bmatrix}1\\0\end{bmatrix}}{\begin{bmatrix}1,0\end{bmatrix}\begin{bmatrix}1\\0\end{bmatrix}}=2$，此时 $x_0=D^{-1}z_0=D^{-1}Q\begin{bmatrix}1\\0\end{bmatrix}=\begin{bmatrix}0\\[1mm]\dfrac{1}{\sqrt{2}}\end{bmatrix}$，
于是 $\max\limits_{x\ne0}\dfrac{f(x)}{g(x)}=2$.
\end{fcard}

% -----------------------------------------------------------
% PDF p120（印刷页 109）—— 全页仅一个【注】框（印刷正文）
% -----------------------------------------------------------
\begin{fnote}
【注】本题是作者新命制的题目，事实上，$\dfrac{x^{\mathrm{T}}Ax}{x^{\mathrm{T}}Bx}$ 比 $\dfrac{x^{\mathrm{T}}Ax}{x^{\mathrm{T}}x}$ 使用范围更为广泛，后者已经在真题中出现过，前者尚未出现，值得注意. 细心的考生应已发现，$\dfrac{x^{\mathrm{T}}Ax}{x^{\mathrm{T}}Bx}$ 依然要转化到 $\dfrac{z^{\mathrm{T}}Mz}{z^{\mathrm{T}}z}$ 上.
\end{fnote}
% 说明：上框为原稿印刷正文。按 AGENTS.md §7.1 决定 2「不保留 OPD 方法论」，
%       已删去末尾的方法论说教（原稿：「这是"三向解题法"中 $D_{23}$（化归经典形式）的
%       再一次具体体现……注意，是意识，意识为先，有了主动的化归意识，办法自然就来了」），
%       **保留全部考点信息**（新命制题型、适用范围更广、后者已考过、仍需化为 $\frac{z^TMz}{z^Tz}$）。

% -----------------------------------------------------------
% PDF p121（印刷页 110）—— 全书最后一页，整页空白，无正文/习题/附录
% -----------------------------------------------------------
% PDF p121（印刷页 110）为全书最后一页，整页空白（无任何文字、无页脚以外的装饰），按规范不排版。
